
\phantomsection\addcontentsline{toc}{section}{第 VII 报告 --- 上同调与 Chern 类}
\section*{第 VII 报告}

\begin{center}
\textbf{若干经典概形的上同调}\\[0.5em]
\textbf{与}\\[0.5em]
\textbf{Chern 类的上同调理论}\\[1em]
作者\\[0.5em]
J.-P. Jouanolou
\end{center}

全文固定整数 $\nu>0$, 并令
\[
A=\mathbb Z/\nu\mathbb Z;
\]
假定所考虑的一切概形的剩余特征都与整数 $\nu$ 互素. 对这样的概形 $X$, 用 $\mu_X$ 表示 $\nu$ 次单位根层, 用 $A^{\bullet}(X)$ 表示上同调环
\[
\bigoplus_i H^{2i}(X,\mu_X^{\otimes i}),
\]
其乘法由杯积给出.

\subsection*{1. 向量丛}

设 $S$ 为概形, $E$ 为局部自由 $\mathcal O_S$-模, 这里约定它为有限型. 令 $X=\V(E)$, 并令 $U=\V(E)^*$ 为零截面的补集; 用
\[
p:X\to S,
\qquad
q:U\to S
\]
表示典范投影.

\begin{statement}{命题 1.1}
设 $L$ 为 $A$-模.
\begin{enumerate}[label=(\roman*)]
\item
\begin{enumerate}[label=\alph*)]
\item 对每个整数 $j>0$, 有
\[
R^j p_*(p^*L)=0.
\]
\item 伴随态射
\[
L\longrightarrow p_*p^*(L)
\]
是同构.
\end{enumerate}
\item 假设 $E$ 的秩恒为 $r$.
\begin{enumerate}[label=\alph*)]
\item
\[
R_!^i p(p^*L)=0\quad\text{若 }i\ne 2r.
\]
\item 迹态射 (SGA 4 XVIII)
\[
R_!^{2r}p(p^*L)\longrightarrow L\otimes\mu_S^{\otimes -r}
\]
是同构.
\end{enumerate}
\end{enumerate}
\end{statement}

\emph{证明.} 对 Zariski 拓扑作局部化, (i) 立即由 (SGA 4 XV 2.2) 得到. 断言 (ii) 的实质已包含在 (SGA 4 XVIII) 中.

\begin{statement}{推论 1.2}
设 $L$ 为 $A_S$-模. 对每个整数 $i$, 逆像态射
\[
H^i(S,L)\longrightarrow H^i(X,p^*L)
\]
是同构.
\end{statement}

\emph{证明.} 由 (1.1), 伴随态射 $L\to Rp_*(p^*L)$ 是同构. 又因为 $H^*(X,p^*L)\simeq H^*(S,Rp_*p^*(L))$, 断言随即成立.

\begin{statement}{命题 1.3}
假设 $E$ 的秩恒为 $r$. 对每个 $A_S$-模 $L$, 下列性质成立.
\begin{enumerate}[label=(\roman*)]
\item
\begin{enumerate}[label=\alph*)]
\item
\[
R^i q_*(q^*L)=0\quad\text{若 }i\ne 0,2r-1.
\]
\item 伴随态射
\[
L\longrightarrow q_*q^*(L)
\]
是同构.
\item 若 $Y$ 表示 $X$ 中 $U$ 的补集, 则态射
\[
R^{2r-1}q_*(q^*L)\longrightarrow L\otimes\mu_S^{\otimes -r}
\]
是同构; 此态射是连接态射
\[
R^{2r-1}q_*(q^*L)\longrightarrow R_Y^{2r}p_*(p^*L)
\tag{1}
\]
与纯性同构 (SGA 4 XVIII)
\[
R_Y^{2r}p_*(p^*L)\longrightarrow L\otimes\mu_S^{\otimes -r}
\]
的复合.
\end{enumerate}
\item
\begin{enumerate}[label=\alph*)]
\item
\[
R_!^i q(q^*L)=0\quad\text{若 }i\ne 1,2r.
\]
\item 连接态射
\[
L\longrightarrow R_!^1 q(q^*L)
\]
是同构; 它来自下列无限正合列, 其中 $f:Y\to S$ 为典范投影:
\[
\cdots\longrightarrow R_!^i q(q^*L)\longrightarrow R_!^i p(p^*L)
\longrightarrow R^i f_*(f^*L)\longrightarrow R_!^{i+1}q(q^*L)\longrightarrow\cdots
\]

\item 迹态射
\[
R_!^{2r}q(q^*L)\longrightarrow L\otimes\mu_S^{\otimes -r}
\]
是同构.
\end{enumerate}
\end{enumerate}
\end{statement}

\emph{证明.} 设 $\alpha:U\to X$ 与 $\beta:Y\to X$ 为典范浸入. 已知 (SGA 4 XVI 3.7), 令 $F=p^*(L)$, 则有
\[
F\xrightarrow{\sim}\alpha_*\alpha^*(F)
\quad\text{且}\quad
(R^j\alpha_*)\alpha^*(F)=0\quad\text{若 }j\ne0,2r-1.
\]
此外, 对 $j\ge1$, 层 $(R^j\alpha_*)\alpha^*F$ 的支集包含于 $Y$, 因而对函子 $p_*$ 无上同调. 由这些事实, 利用 Leray 谱序列
\[
R^i p_*\bigl(R^j\alpha_*\bigr)(q^*L)\Longrightarrow R^n q_*(q^*L).
\]
便得到 (i) a) 与 b). 断言 (i) c) 由 (1.1 (i)) 以及无限正合列
\[
\cdots\longrightarrow R_Y^i p_*(p^*L)\longrightarrow R^i p_*(p^*L)\longrightarrow R^i q_*(q^*L)
\longrightarrow R_Y^{i+1}p_*(p^*L)\longrightarrow\cdots.
\]
得到. 层 $R_Y^{2r}p_*(p^*L)$ 伴随于预层
\[
U\longmapsto H^{2r}_{Y\times_S U}(X\times_S U,p^*L).
\]

下面证明 (ii). 断言 c) 来自 $q$ 光滑且几何纤维连通、维数为 $r$ 这一事实. b) 中的正合列是对典范正合列
\[
0\longrightarrow \alpha_!(q^*L)\longrightarrow p^*L\longrightarrow \beta_!(\beta^*p^*L)\longrightarrow 0.
\]
施以函子 $R_!p$ 而得到的. 利用 (1.1 (i)), 从这一无限正合列不难得到 (ii) a) 与 b).

\begin{statement}{注 1.4}
当 $S=\operatorname{Spec}(\mathbb C)$ 时, (1.3) 表明 $\mathbb C^r$ 中原点的补集, 与作为 $\mathbb C^r-0$ 的形变收缩核的 $2r-1$ 维实球面有相同的上同调.
\end{statement}

\begin{statement}{推论 1.5 (Gysin 正合列)}
Leray 谱序列
\[
H^i(S,R^j q_*(q^*L))\Rightarrow H^{i+j}(U,q^*L)
\]
给出称为 Gysin 正合列的无限正合列
\[
\cdots\longrightarrow H^i(S,L)\longrightarrow H^i(U,q^*L)\longrightarrow
H^{i-2r+1}(S,L\otimes\mu_S^{\otimes -r})\longrightarrow H^{i+1}(S,L)\longrightarrow\cdots.
\]
\end{statement}

\emph{证明.} 应用二行引理.

\subsection*{2. 射影概形}

下面采用的叙述方式基于 L. Illusie 的建议, 所用方法接近 P. Deligne 为证明某些谱序列退化而使用的方法 [6].

\subsubsection*{2.1}
设 $X,Y$ 为两个概形, $f:X\to Y$ 为态射. 我们将在 $D(Y)$ 中定义箭头
\[
Rf_*(A_X)\otimes^L Rf_*(A_X)\longrightarrow Rf_*(A_X)
\tag{2.1.1}
\]
它具有通常的结合性和交换性, 从而赋予 $Rf_*(A_X)$ 广义代数结构. 为此, 记
\[
C^*(X,A_X)
\]
为 $A_X$ 的 Godement 消解 (SGA 4 XVII). 如同 (TF 6.6) 中那样, 可以定义消解的态射
\[
C^*(X,A_X)\otimes C^*(X,A_X)\longrightarrow C^*(X,A_X\otimes A_X).
\]
将它与由 $A_X$ 的乘法导出的态射
\[
C^*(X,A_X\otimes A_X)\longrightarrow C^*(X,A_X)
\]
复合, 最终得到态射
\[
C^*(X,A_X)\otimes C^*(X,A_X)\longrightarrow C^*(X,A_X).
\tag{2.1.2}
\]
施以函子 $f_*$, 便不难得到 $A_Y$-模复形的态射
\[
f_*C^*(X,A_X)\otimes f_*C^*(X,A_X)\longrightarrow f_*C^*(X,A_X),
\]
再对左边表达式的各分量取平坦消解, 便得到 (2.1.1).

\subsubsection*{2.2}
设 $S$ 为概形, $E$ 为秩 $r+1$ 的局部自由 $\mathcal O_S$-模. 用
\[
P=P_S(E)\xrightarrow{p}S
\]
表示相应的射影丛. 已知典范 $\mathcal O_P$-模 $\mathcal O_P(1)$ 通过 Kummer 正合列确定元素
\[
\xi\in H^2(P,\mu),
\]
用 $\eta$ 表示 $\xi$ 在典范态射
\[
H^2(P,\mu)\longrightarrow H^0(S,R^2p_*(\mu))
\]
下的像, 此态射来自 Leray 谱序列
\[
H^i(S,R^jp_*(\mu))\Rightarrow H^{i+j}(P,\mu).
\]
将类 $\xi$ 等同于 $D(P)$ 中的箭头 $A_P\to\mu_P[2]$, 转置后得到 $\mu_P^{\otimes -1}[-2]\to A_P$, 再由通常的伴随得到箭头
\[
c:\mu_S^{\otimes -1}[-2]\longrightarrow Rp_*(A_P),
\]
它在 $2$ 维上同调上与 $\eta$ 一致.

通过张量积, 箭头 $c$ 借助 (2.2.1) 对每个整数 $i\ge1$ 确定箭头
\[
c_i:\mu_S^{\otimes -i}[-2i]\longrightarrow Rp_*(A_P),
\]
并约定用 $c_0$ 表示伴随箭头
\[
A_S\longrightarrow Rp_*(A_P).
\]
于是得到 $D(S)$ 中的箭头
\[
\gamma=\bigoplus_{i=0}^{r}c_i:
\bigoplus_{i=0}^{r}\mu_S^{\otimes -i}[-2i]\longrightarrow Rp_*(A_P)
\]
.

\begin{statement}{定理 2.2.1}
箭头 $\gamma$ 是同构.
\end{statement}

箭头 $\gamma$ 一旦定义, (2.2.1) 便等价于下列推论.

\begin{statement}{命题 2.2.2}
在上述假设下, 有:
\begin{enumerate}[label=\alph*)]
\item
\[
R^{2i+1}p_*(A_P)=0\quad\text{对每个 }i\in\mathbb N.
\]
\item 赋予 $\bigoplus_i R^{2i}p_*(\mu_P^{\otimes i})$ 由杯积导出的 $A_S$-代数结构后, 由 $\eta$ 确定的 $A_S$-代数态射
\[
A_S[T]\longrightarrow \bigoplus_i R^{2i}p_*(\mu_P^{\otimes i})
\]
的核是理想 $(T^{r+1})$, 因而给出 $A_S$-代数同构
\[
\theta:A_S[T]/(T^{r+1})\xrightarrow{\sim}\bigoplus_i R^{2i}p_*(\mu_P^{\otimes i}).
\]
\item 态射
\[
A_S\longrightarrow R^{2r}p_*(\mu^{\otimes r})
\]
是迹态射 (SGA 4 XVIII) 的逆; 它由 $\theta$ 与 $A_S[T]/(T^{r+1})$ 上乘以 $T^r$ 的映射在 $A_S$ 上的限制复合而得:
\[
R^{2r}p_*(\mu^{\otimes r})\longrightarrow A_S;
\]
换言之,
\[
\operatorname{Tr}(\eta^r)=1.
\]
\end{enumerate}
\end{statement}

\emph{证明.} 由于这些资料在基变换下稳定, 固有基变换定理 (SGA 4 XII 5.1) 以及迹态射与基变换的相容性 (SGA 4 XVIII), 使我们只需在 $P$ 的各纤维上验证相应断言. 换言之, 可以假设 $S$ 是可分闭域 $k$ 的谱, $p$ 是典范态射 $\mathbb P^r_k\to k$. 在这种情形下, 证明使用以下引理.

\begin{statement}{引理 2.2.3}
设 $k$ 为可分闭域, $p:\mathbb P_k^r\to k$ 为典范投影. 则:
\begin{enumerate}[label=\alph*)]
\item
\[
R^{2i+1}p_*(A_P)=0\quad\text{对每个 }i.
\]
\item
\[
R^{2i}p_*(A_P)\simeq A_k\quad\text{对 }i\in[0,r],
\qquad
R^{2i}p_*(A_P)=0\quad\text{若 }i>r.
\]
\end{enumerate}
\end{statement}

本引理在 $r=0$ 时显然成立, 对 $r$ 归纳即可证明. 假定它对 $r-1$ 成立, 来证明对 $r$ 成立. 射影空间 $\mathbb P_k^{r-1}$ 作为闭子空间嵌入 $\mathbb P_k^r$, 补集为开集 $\mathbb E_k^r$, 因此有固有支集上同调的正合列
\[
\cdots\longrightarrow H_!^i(\mathbb E_k^r)\longrightarrow H^i(\mathbb P_k^r)
\longrightarrow H^i(\mathbb P_k^{r-1})\longrightarrow H_!^{i+1}(\mathbb E_k^r)\longrightarrow\cdots.
\]
由 $\mathbb E_k^r$ 的固有支集上同调的计算 (1.1), 当 $i\ne(2r,2r-1)$ 时, 逆像映射
\[
H^i(\mathbb P_k^r)\longrightarrow H^i(\mathbb P_k^{r-1})
\]
是同构, 从而此时断言成立. 若 $i=2r-1$, 同样可见 $H^{2r-1}(\mathbb P_k^r)$ 包含于 $H^{2r-1}(\mathbb P_k^{r-1})$, 因而由归纳假设为零. 最后, 若 $i=2r$, 关系
\[
H^{2r-1}(\mathbb P_k^{r-1})=H^{2r}(\mathbb P_k^{r-1})=0
\]
表明
\[
H_!^{2r}(\mathbb E_k^r)\xrightarrow{\sim}H^{2r}(\mathbb P_k^r),
\]
从而由 (1.1) 得到断言.

此引理蕴含 2.2.1 的 a). 现在证明 $\theta$ 是同构. 态射 $\theta$ 可以解释为环态射
\[
A[T]/(T^{r+1})\longrightarrow\bigoplus_i H^{2i}(P,\mu^{\otimes i}).
\]


因此, 根据 (2.2.3), 只需证明对每个 $i\in[0,r]$, $\xi^i$ 是 $A$-模 $H^{2i}(P,\mu^{\otimes i})$ 的自由元素. 又只需对 $\xi^r$ 验证, 而这将直接由 2.2.2 c) 得到.

现在证明 c). 实际上, 在这一情形这就是迹态射的定义; 无论如何, 它表达了迹态射与下图中基本类的相容性:
\[
\begin{tikzcd}[column sep=large,row sep=large]
x_0 \arrow[r,hook,"\mathrm{inj}"] \arrow[dr] & \mathbb P^r_k \arrow[d,"p"]\\
& k ,
\end{tikzcd}
\]
其中 $\mathrm{inj}:x_0\to \mathbb P^r_k$ 是 $\mathbb P^r_k$ 的一个闭点的嵌入.

\begin{statement}{推论 2.2.4}
设 $L$ 为 $A_S$-模复形. 箭头
\[
\mathrm R\Gamma(S,\gamma\lten_A\mathrm{id}(L)):
\bigoplus_{i=0}^r \mathrm R\Gamma(S,L\otimes\mu_S^{\otimes -i})[-2i]
\longrightarrow
\mathrm R\Gamma(P,p^*(L))
\]
是同构.
\end{statement}

事实上, 只需对同构 $\gamma\lten_A\mathrm{id}_L$ 施以函子 $\mathrm R\Gamma_S$.

\subsubsection*{2.2.5}
为陈述下一推论, 对每个概形 $X$ 记
\[
A^*(X)=\bigoplus_i H^{2i}(X,\mu_X^{\otimes i})
\]
以下始终使用这一记号.

\begin{statement}{推论 2.2.6}
逆像同态
\[
p^*:A^*(S)=\bigoplus_i H^{2i}(S,\mu_S^{\otimes i})
\longrightarrow
\bigoplus_i H^{2i}(P,\mu_P^{\otimes i})=A^*(P)
\]
是含幺环的单射态射, 且元素 $1,\xi,\ldots,\xi^r$ 构成 $A^*(S)$-模 $A^*(P)$ 的一组基. 特别地, $A^*(P)$ 是自由 $A^*(S)$-模, 其秩与 $E$ 的秩相同.
\end{statement}

\section*{3. Chern 类}

设 $S$ 为概形, $E$ 为秩恒为 $r$ 的局部自由 $\mathcal O_S$-模. 用 $\check E$ 表示 $E$ 的对偶, 令 $P=P_S(\check E)$. 已经看到 (2.2.6), 若以 $\xi$ 表示 $\mathcal O_P(1)$ 在 $H^2(P,\mu)$ 中的类, 则 $A^*(S)$-模 $A^*(P)$ 是自由的, 基为 $1,\xi,\ldots,\xi^{r-1}$. 特别地, 根据齐次性, 立即可见存在唯一形如
\[
\xi^r+c_1\xi^{r-1}+\cdots+c_i\xi^{r-i}+\cdots+c_r=0,
\tag{3.1}
\]
的关系, 其中对每个 $i\in[0,r]$, 有 $c_i\in H^{2i}(S,\mu^{\otimes i})$.

现在设 $E$ 为任意局部自由 $\mathcal O_S$-模. 对每个整数 $r$, $S$ 中使 $E$ 的秩等于 $r$ 的点组成开集 $S_r$. 于是得到 $S$ 分解成不交开集的分拆
\[
S=\coprod_r S_r
\]
特别地, 对每个 $A_S$-模 $M$ 和每个整数 $i$, 各逆像态射的乘积
\[
H^i(S,M)\longrightarrow \prod_r H^i(S_r,M/S_r)
\]
是双射.

上述讨论说明以下定义是合理的.

\begin{statement}{定义 3.2}
设 $S$ 为概形, $E$ 为局部自由 $\mathcal O_S$-模. 对每个整数 $i\ge0$, $E$ 的第 $i$ 个 Chern 类记为 $c_i(E)$, 它是 $H^{2i}(S,\mu^{\otimes i})$ 中的元素, 对每个整数 $r\ge0$, 其在 $E$ 的秩等于 $r$ 的点构成的开集 $S_r$ 上的限制, 等于关系 3.1 定义的元素 $c_i$. 若 $E$ 的秩有界, 则 $E$ 的总 Chern 类是 $A^*(S)$ 中的元素
\[
c(E)=\sum_i c_i(E)
\]
一般情形下, 此记号表示下列积中的相应元素:
\[
\widehat A^*(S)=\prod_i A^i(S).
\]
\end{statement}

\subsubsection*{3.3}
显然, 对所有 $S,E$, 有 $c_0(E)=1$; 若 $E$ 的秩以 $d$ 为界, 则当 $i>d$ 时有 $c_i(E)=0$.

\begin{statement}{命题 3.4}
刚才定义的 Chern 类具有下列性质.
\begin{enumerate}[label=(\roman*)]
\item \emph{归一化.} 设 $L$ 为可逆 $\mathcal O_S$-模. 若以 $[L]$ 表示 $L$ 的类在连接态射
\[
H^1(S,\mathbb G_m)\longrightarrow H^2(S,\mu),
\]
下的像, 则
\[
c(L)=1+[L].
\]
\item \emph{基变换下的稳定性.} 设 $f:S'\to S$ 为概形态射, $E$ 为局部自由 $\mathcal O_S$-模. 对每个整数 $i$,
\[
f^*(c_i(E))=c_i(f^*E).
\]
\item \emph{可加性.} 设
\[
0\to E'\to E\to E''\to0
\]
为局部自由 $\mathcal O_S$-模的正合列. 则对每个整数 $n$, 有
\[
c_n(E)=\sum_{i+j=n}c_i(E')c_j(E'').
\]
\end{enumerate}
\end{statement}

开始证明前, 先抽出以下推论.

\begin{statement}{推论 3.5}
对每个局部自由 $\mathcal O_S$-模的正合列 $0\to E'\to E\to E''\to0$, 有
\[
c(E)=c(E')c(E'')
\]
等式在 $\widehat A^*(S)$ 中成立.
\end{statement}

\emph{3.4 的证明.} 遵循 A. Grothendieck 的文章: \emph{La théorie des classes de Chern}, Bull. S.M.F. (1958), pp. 137--154.

(i) 设 $E$ 为局部自由 $\mathcal O_S$-模; 令 $P=P_S(\check E)$, $p:P\to S$ 为典范投影. 回顾存在典范满同态
\[
p^*(\check E)\longrightarrow \mathcal O_P(1).
\]
特别地, 若 $E$ 可逆, 并按 (i) 的记号令 $E=L$, 则 $P=S$, 且有同构
\[
L\xrightarrow{\sim}\mathcal O_P(1).
\]
因此 $\xi+[L]=0$, 断言由 Chern 类的定义得到.

(ii) 验证是直接的, 留给读者.

(iii) 将 $S$ 分拆成开集, 使有关模在每一开集上的秩恒定, 可以假设 $E,E',E''$ 的秩分别恒为 $r,r',r''$. 假设已证明以下引理, 先说明如何由它得到断言.

\begin{statement}{引理 3.6}
设 $E$ 为局部自由 $\mathcal O_S$-模, 配有有限滤过
\[
0=E_0\subset E_1\subset\cdots\subset E_i\subset\cdots\subset E_r=E
\]
且相邻商 $L_i=E_i/E_{i-1}$ $(1\le i\le r)$ 均可逆. 则
\[
c(E)=(1+[L_1])(1+[L_2])\cdots(1+[L_r]).
\]
\end{statement}

设 $h:Z\to S$ 为 $E'$ 和 $E''$ 的旗概形的纤维积. 由旗概形的定义, 显然 $h^*(E')$ 和 $h^*(E'')$ 分别具有相邻商为可逆模 $(L'_i)$ $(1\le i\le r')$ 与 $(L''_j)$ $(1\le j\le r'')$ 的有限滤过. 因而 $h^*(E)$ 具有以各 $L'_i$ 和 $L''_j$ 为相邻商的滤过. 应用引理 3.6 得到
\[
c(h^*(E))=
\prod_{1\le i\le r'}c(L'_i)\prod_{1\le j\le r''}c(L''_j)
=c(h^*(E'))c(h^*(E'')),
\]
再由 (ii), 得到
\[
h^*(c(E))=h^*(c(E')c(E'')).
\tag{3.5.1}
\]
而 $h$ 是形如
\[
P_T(F)\longrightarrow T
\qquad (F\text{ 为局部自由 }\mathcal O_T\text{-模}),
\]
的态射的复合, 故由 (2.2.5), 逆像同态 $h^*:A^*(S)\to A^*(Z)$ 为单射. 结合 (3.5.1) 即得所需断言.

下面证明引理 3.6. 由前述论证, 只需证明相应的关系
\[
c(p^*E)=\prod_i c(p^*(L_i)),
\]
其中回顾 $p:P=P_S(\check E)\to S$ 为典范投影. 对每个 $i\in[1,r]$, 单态射 $E_i\to E$ 通过对偶确定满态射 $\check E\to\check E_i$; 令 $Y_i=P_S(\check E_i)$, 得到 $P$ 的一列闭子概形:
\[
S\simeq Y_1\subset Y_2\subset\cdots\subset Y_i\subset\cdots\subset Y_r=P.
\]
对每对满足 $1\le i\le j\le r$ 的整数 $(i,j)$, 设
\[
\alpha_{ji}:Y_i\to Y_j
\]
为典范闭浸入, 并令 $\alpha_i=\alpha_{r,i}$. 对每个整数 $i\in[1,r-1]$, 由下面显然的引理 (3.6.1), $Y_i$ 在 $Y_{i+1}$ 中由同构于
\[
\alpha_{i+1}^*(p^*(\check L_{i+1})\otimes\mathcal O_P(-1)).
\]
的可逆理想定义.
\begin{statement}{引理 3.6.1}
设 $0\to L\to E\to F\to0$ 为局部自由 $\mathcal O_S$-模的正合列. 则下列序列
\[
0\to L\otimes_S S(E)\xrightarrow{u}S(E)\xrightarrow{v}S(F)\to0,
\]
是分次 $\mathcal O_S$-模的正合列, 其中 $v$ 为显然的态射, $u$ 为对应于 $L$ 嵌入 $E$ 的一次态射.
\end{statement}

此外, $Y_1$ 等同于 $S$, 并且有同构
\[
\alpha_1^*(p^*(\check L_1))\simeq\mathcal O_{Y_1}(1),
\]
换言之,
\[
\alpha_1^*(p^*(\check L_1)\otimes\mathcal O_P(-1))\simeq \mathcal O_{Y_1}.
\tag{3.6.2}
\]
为得出结论, 我们使用与光滑 $S$-概形的闭浸入 $u:Z\to X$ 相关联的 Gysin 同态 (SGA 4 XVIII) 的下列两个性质; 该浸入由 $\mathcal O_X$ 的理想 $J$ 定义. 它满足投影公式
\[
u_*(x\,u^*(y))=u_*(x)y,
\tag{3.6.3}
\]
且当理想 $J$ 可逆时满足关系
\[
u_*(1_Z)=-[J]=[\check J].
\tag{3.6.4}
\]
事实上, 由此将对整数 $i$ 作递增归纳, 证明
\[
(\alpha_i)_*(\alpha_i)^*\Bigl(\prod_{j\le i}[p^*(\check L_j)\otimes\mathcal O_P(-1)]\Bigr)=0.
\tag{3.6.5}
\]
当 $i=r$ 时, 断言 (3.6.5) 表明
\[
\prod_i[p^*(\check L_i)\otimes\mathcal O_P(-1)]=0,
\]
或等价地,
\[
\prod_i(\xi+c_1(p^*(\check L_i)))=0,
\]
这立即给出 (3.6). 当 $i=1$ 时, 断言 (3.6.5) 由 (3.6.2) 得到. 假设对 $i$ 成立, 来证明对 $i+1$ 成立. 为简化记号, 令
\[
N=\prod_{j\le i}[p^*(\check L_j)\otimes\mathcal O_P(-1)].
\]
对 $i$ 的等式 (3.6.5) 也可写为
\[
(\alpha_{i+1})_*(\alpha_{i+1,i})_*(\alpha_{i+1,i})^*(\alpha_{i+1})^*(N)=0,
\]
于是, 由 $\alpha_{i+1,i}$ 的投影公式,
\[
(\alpha_{i+1})_*\bigl[(\alpha_{i+1})^*(N)(\alpha_{i+1})^*(p^*(\check L_{i+1})\otimes\mathcal O_P(-1))\bigr]=0,
\]
这正是所需结果.

\begin{statement}{推论 3.7}
\begin{enumerate}[label=(\roman*)]
\item 对每个局部自由 $\mathcal O_S$-模 $E$ 和每个整数 $i$,
\[
c_i(\check E)=(-1)^i c_i(E).
\]
\item 若 $E$ 为秩 $r$ 的局部自由 $\mathcal O_S$-模, 则
\[
c_1(E)=c_1(\Lambda^rE).
\]
\item 更一般地, 张量积或外幂的 Chern 类, 可用 (SGA 6 V 6.1) 中明确给出的泛多项式表达.
\end{enumerate}
\end{statement}

\emph{证明.} 例如证明 (ii), 其余断言同理. 利用 $E$ 的旗概形, 可归约到 $E$ 具有相邻商为可逆模 $L_i$ 的滤过的情形. 此时断言由熟知的同构
\[
\Lambda^r(E)\simeq\bigotimes_i L_i.
\]
得到.
\subsubsection*{3.8}
给定在 $\operatorname{Spec}(\mathbb C)$ 上局部有限型的概形 $X$, 其有理点集 $X(\mathbb C)$ 典范地带有局部紧拓扑空间结构. 因而代数拓扑给出 Chern 类理论 [3], 给 $X(\mathbb C)$ 上每个连续复向量丛 $V$ 配以类 $c_i(V)\in H^{2i}(X(\mathbb C),\mathbb Z)$. 现在设 $\mathcal O_\mathbb C$ 为 $X(\mathbb C)$ 上连续复值函数层, $\mathcal E$ 为局部自由有限型 $\mathcal O_\mathbb C$-模. 定义 Chern 类
\[
c_i(\mathcal E)\in H^{2i}(X(\mathbb C),\mathbb Z),
\]
即连续截面层为 $\mathcal E$ 的复向量丛的 Chern 类. 最后, 利用显然的映射 $\mathbb Z\to\mathbb Z/\nu\mathbb Z$, 得到 Chern 类 (不会混淆时仍用相同记号)
\[
c_i(\mathcal E)\in H^{2i}(X(\mathbb C),\mathbb Z/\nu\mathbb Z).
\tag{3.8.1}
\]
另一方面, 在 $X_{\et}$ 上配以结构层 $\mathcal O_a$, 在 $X(\mathbb C)$ 上配以环 $\mathcal O_\mathbb C$ 后, 有赋环拓扑斯的态射 (SGA 4 XI)
\[
\varepsilon:X(\mathbb C)\to X_{\et},
\]
其环的结构态射
\[
\varepsilon^*(\mathcal O_a)\to\mathcal O_\mathbb C
\tag{3.8.2}
\]
来自每个代数函数均连续这一事实. 态射 $\varepsilon$ 诱导阿贝尔层的态射
\[
\varepsilon^*(\mu_{X_{\et}})\to \mu_{X(\mathbb C)},
\tag{3.8.3}
\]
其中 $\mu$ 在两种背景下都表示 $\nu$ 次单位根层. 最后, 通过同构
\[
(\mathbb Z/\nu\mathbb Z)_{X(\mathbb C)}
 \xrightarrow{\sim}\mu_{X(\mathbb C)}
\tag{3.8.4}
\]
将 $(\mathbb Z/\nu\mathbb Z)_{X(\mathbb C)}$ 与 $\mu_{X(\mathbb C)}$ 典范地等同, 此同构把截面 $1$ 送到截面 $e^{2i\pi/\nu}$.

\begin{statement}{命题 3.8.5}
设 $E$ 为局部自由 $\mathcal O_X$-模. 对每个整数 $i\ge0$, 设
\[
\varepsilon^*:H^{2i}(X_{\et},\mu_X^{\otimes i})
\to H^{2i}(X(\mathbb C),\mathbb Z/\nu\mathbb Z)
\]
为从 3.8.3 和 3.8.4 导出的逆像态射. 则
\[
c_i(\varepsilon^*E)=\varepsilon^*c_i(E).
\]
\end{statement}

\emph{证明.} 只需在 $E$ 可逆时验证. 此时作通常的等同后, $\varepsilon^*(E)$ 是 $E$ 在典范映射
\[
\varepsilon^*:H^1(X_{\et},\mathbb G_m)\to H^1(X(\mathbb C),\mathcal O_\mathbb C^*),
\tag{3.8.6}
\]
下的像; 我们要证明下图交换:
\[
\begin{tikzcd}[column sep=large,row sep=large]
H^1(X,\mathbb G_m) \arrow[r,"\varepsilon^*"] \arrow[d,"\partial"'] &
H^1(X(\mathbb C),\mathcal O_\mathbb C^*) \arrow[d,"c_1"]\\
H^2(X,\mu_X) \arrow[r,"\varepsilon^*"'] &
H^2(X(\mathbb C),\mathbb Z/\nu\mathbb Z)
\end{tikzcd}
\tag{3.8.7}
\]
其中 $\partial$ 是来自典范正合列
\[
0\to \mu_X\to\mathbb G_m\xrightarrow{\nu}\mathbb G_m\to0
\tag{3.8.8}
\]
的连接态射, $c_1$ 是第一 Chern 类 (3.8.1). 事实上, 利用“第一 Chern 类”态射
\[
H^1(X(\mathbb C),\mathcal O_\mathbb C^*)\to H^2(X,\mathbb Z)
\]
是典范正合列
\[
0\to\mathbb Z\to\mathcal O_\mathbb C\xrightarrow{\exp}\mathcal O_\mathbb C^*\to0,
\qquad
t\mapsto 2i\pi t,
\]
的连接态射这一事实, 容易看出 (3.8.7) 的态射 $c_1$ 正是正合列
\[
0\to\mathbb Z/\nu\mathbb Z\to\mathcal O_\mathbb C^*\xrightarrow{\nu}\mathcal O_\mathbb C^*\to0,
\qquad
t\mapsto e^{2i\pi t/\nu},
\]
所关联的连接态射, 而由 (3.8.4), 它与 (3.8.8) 对应. 断言立即成立.

\subsubsection*{3.9}
设 $X$ 为概形, $E$ 为局部自由 $\mathcal O_X$-模. 直接取极限, 便可对每个与 $X$ 的剩余特征互素的素数 $\ell$, 定义 $\ell$ 进 Chern 类
\[
c_i(E)\in H^{2i}(X,\mathbb Z_\ell(i))
\stackrel{\mathrm{def}}{=}\varprojlim_\nu H^{2i}(X,\mu_{\ell^\nu}^{\otimes i}),
\]
它们具有 3.4 和 3.7 中列出的通常性质, 且在 3.8 的背景下, 显然在自然意义下与拓扑定义的 Chern 类相容.

与 $\mathbb Q_\ell$ 作张量积, 得到模挠的 $\ell$ 进 Chern 类
\[
c_i(E)\in H^{2i}(X,\mathbb Q_\ell(i))
=H^{2i}(X,\mathbb Z_\ell(i))\otimes_{\mathbb Z_\ell}\mathbb Q_\ell,
\]
它们显然也满足相同性质.

现在假设 $X$ 是域 $k$ 上拟射影光滑概形, 这里用 $C^*(X)$ 表示其 Chow 环 (Sém. Chevalley: Anneaux de Chow et applications III 4). 回顾在 IV 中已定义分次环态射
\[
\varphi_X:C^*(X)\to A^*(X)
\]
它与正像及逆像态射相容. 此外, 典范同构
\[
\varepsilon_X:\Pic(X)\xrightarrow{\sim} C^1(X)
\]
将每个可逆 $\mathcal O_X$-模 $L$ 送到 $L$ 的一个在 $X$ 的任一连通分支上均不为零的有理截面的除子类; 沿用 Grothendieck 上述文章中的方法 (该文假设 $k$ 代数闭), 便得到取值于 $C^*(X)$ 的 Chern 类理论. 复合
\[
\varphi_X\circ\varepsilon_X:\Pic(X)\longrightarrow H^2(X,\mathbb Z_\ell(1))
\]
正是第一 $\ell$ 进 Chern 类. 事实上, 只需对形如 $\check J$ 的 $\Pic(X)$ 元素验证, 其中 $J$ 为 $\mathcal O_X$ 的可逆理想; 此时由 (3.6.4) 即得. 根据 Chern 类理论的唯一性, 态射 $\varphi_X$ 将 Chern 类送到 Chern 类.

\section*{4. 自交公式及应用}

设 $S$ 为概形, $X$ 为 $S$-概形, $Y$ 为 $X$ 的闭子概形, 它在 $S$ 上光滑, 且处处余维为 $d$. 最后假设 $X$ 在 $Y$ 的某个邻域内在 $S$ 上光滑. 设 $i:Y\to X$ 为典范闭浸入, $J$ 为定义 $Y$ 的 $\mathcal O_X$-理想. 回顾 (SGA $4^{1/2}$ Cycle), 此时有 Gysin 同构
\[
i_*:H^{2p}(Y,\mu_Y^{\otimes p})
\to H_Y^{2p+2d}(X,\mu_X^{\otimes(p+d)}),
\]
即与基本类的杯积. 不会混淆时, 仍用相同符号表示由这些同构忘掉支集而得的 Gysin 态射
\[
i_*:H^{2p}(Y,\mu_Y^{\otimes p})
\to H^{2p+2d}(X,\mu_X^{\otimes(p+d)})
\]
.

\begin{statement}{定理 4.1}
在上述情形下, 若
\[
N=J/J^2
\]
表示 $i$ 的法层, 则对每个 $y\in A^*(Y)$, 有
\[
i^*i_*(y)=y\,c_d(\check N).
\]
\end{statement}

开始证明这个定理前, 先说明几个变体, 它们的证明完全相同.

a) 若 $\ell$ 为与 $S$ 的剩余特征互素的素数, 则相应公式在 $H^*(Y,\mathbb Z_\ell(*))$ 中成立.

b) 可以加入支集. 例如, 若 $T$ 为 $Y$ 的闭子集, $y\in\allowbreak  H_T^{2p}(Y,\mu_Y^{\otimes p})$, 则公式 4.1 在 $H_T^{2p+2d}(Y,\mu_Y^{\otimes(p+d)})$ 中仍成立. 当然, 此时 $i_*$ 表示 Gysin 态射
\[
H_T^{2p}(Y,\mu_Y^{\otimes p})
\to H_T^{2p+2d}(X,\mu_X^{\otimes(p+d)}),
\]
它仍是同构 (SGA $4^{1/2}$ Cycle).

c) 最后, 可以将这两种情形结合起来.

由于 4.1 的证明的基本工具是爆破概形理论, 现在回顾若干相关事实. 设
\[
f:X'\to X
\]
为沿 $Y$ 爆破所得的 $X$-概形, 并设
\[
\begin{tikzcd}[column sep=large,row sep=large]
Y' \arrow[r,"j"] \arrow[d,"g"'] & X' \arrow[d,"f"]\\
Y \arrow[r,"i"'] & X
\end{tikzcd}
\]
为由 $f$ 与 $i$ 构成的笛卡尔图. 回顾 (参见 SGA 6 VII 3.1, 这里采用其记号), $j$ 是余维 $1$ 的正则浸入, 由理想
\[
J'\simeq\mathcal O_{X'}(1)
\]
定义, 且 $Y'\simeq P_Y(N)$. 此外, 有 $\mathcal O_{Y'}$-模的正合列
\[
0\to F\to g^*(N)\to \mathcal O_{Y'}(1)\to0.
\]

\begin{statement}{引理 4.2}
对每个 $A_X$-模 $M$, 同态
\[
f^*:H_Y^*(X,M)\to H_{Y'}^*(X',f^*(M))
\]
与
\[
f_*:H_{Y'}^*(X',f^*(M))\to H_Y^*(X,M)
\]
满足关系 $f_*f^*=\mathrm{id}$. 此外, $f_*$ 的次数为 $0$.
\end{statement}

\emph{证明.} 由于 $f$ 是相对维数为零的完全交态射 (SGA 6 VII 1.8), 最后一个断言由 Gysin 同态的定义得到 (SGA $4^{1/2}$ Cycle). 对于第一个断言, $f$ 的投影公式
\[
f_*(x\,f^*(x'))=f_*(x)x',
\]
(其中 $x\in H^*(X,A)$, $x'\in H_Y^*(X,M)$) 将问题归约为证明 $f_*(1)=1$. 只需在 $U=X-Y$ 上验证相应断言; 因为 $f$ 诱导 $U'=X'-Y'$ 到 $U$ 的同构, 这是显然的.

\begin{statement}{注 4.3}
当然, 忘掉支集后得到的 4.2 的相应结论也成立, 证明相同.
\end{statement}

现在证明 4.1. 首先, 由 (SGA $4^{1/2}$ Cycle) 中一个本质上以形式方法证明的一般结果, 有
\[
i^*i_*(y)=y\,i^*i_*(1),
\]
所以只需证明
\[
i^*i_*(1_Y)=c_d(\check N),
\]
这里 $i_*$ 表示 Gysin 同构
\[
H^0(Y,A_Y)\to H_Y^{2d}(X,\mu_X^{\otimes d}).
\]
由于 $f^*i_*1_Y\in H_{Y'}^{2d}(X',\mu_{X'}^{\otimes d})$, 且 Gysin 态射
\[
j_*:H^{2d-2}(Y',\mu_{Y'}^{\otimes(d-1)})
\to H_{Y'}^{2d}(X',\mu_{X'}^{\otimes d})
\]
是同构, 存在
\[
u\in H^{2d-2}(Y',\mu_{Y'}^{\otimes(d-1)})
\]
使得
\[
f^*i_*(1_Y)=j_*(u).
\tag{4.1.1}
\]
对 4.1.1 两边施以 $j^*$, 得到
\[
j^*f^*i_*(1_Y)=j^*j_*(u).
\]
令 $\xi=c_1(\mathcal O_{Y'}(1))$, 由 3.6.4 得到
\[
g^*i^*i_*(1_Y)=-u\,\xi.
\tag{4.1.2}
\]
元素 $u$ 具有唯一的下述分解:
\[
u=a_0+a_1\xi+\cdots+a_{d-1}\xi^{d-1},
\qquad
a_i\in H^{2d-2-2i}(Y,\mu_Y^{\otimes(d-1-i)}).
\tag{4.1.3}
\]
由 $\check N$ 的 Chern 类的定义关系
\[
\xi^d+c_1(\check N)\xi^{d-1}+\cdots+c_{d-1}(\check N)\xi+c_d(\check N)=0,
\tag{4.1.4}
\]
得到等式
\[
u\xi=-a_{d-1}c_d(\check N)+(a_0-a_{d-1}c_{d-1}(\check N))\xi
+\cdots+(a_{d-2}-a_{d-1}c_1(\check N))\xi^{d-1}.
\]
将最后这个表达式与 (4.1.2) 比较, 利用 $g^*$ 的单射性 (2.2.6), 得到
\[
i^*i_*(1_Y)=a_{d-1}c_d(\check N).
\tag{4.1.5}
\]
还需确定 $a_{d-1}$. 首先, 同态 $g_*$ 将次数降低 $2d-2$, 因而
\[
g_*(\xi^i)=0\quad(0\le i\le d-2),
\]
此外, 由 (2.2.2),
\[
g_*(\xi^{d-1})=1.
\]
对 4.1.3 两边施以 $g_*$, 得到
\[
g_*(u)=a_{d-1}.
\tag{4.1.6}
\]
另一方面, 对 4.1.1 两边施以 $f_*$, 由 4.2 得到
\[
i_*(1_Y)=i_*g_*(u).
\tag{4.1.7}
\]
而这里的同态 $i_*$ 是单射, 因此
\[
a_{d-1}=g_*(u)=1_Y,
\]
证明完毕.

\begin{statement}{注 4.4}
a) 上述论证的关键是纯性定理. 一旦获得完全一般的纯性定理, 当 $Y$ 正则且在 $X$ 中纯余维为 $d$, $X$ 卓越且在 $Y$ 的一个邻域内正则时, 命题 4.1 就成立. 无论如何, 在特征 $0$ 时, 这种形式的结果已由 (SGA 4 XIX) 建立.

b) 当 $S$ 为可分闭域的谱, $X,Y$ 为 $S$ 上拟射影光滑概形时, Mumford 已在 $Y$ 的 Chow 环中建立相应断言 (参见 9.2).
\end{statement}

\begin{statement}{推论 4.5}
在 4.1 的假设下, 设 $x,y\in A^*(Y)$. 则
\[
i_*(x)i_*(y)=i_*(xy\,c_d(\check N)).
\]
\end{statement}

\emph{证明.} 态射 $i$ 的投影公式表明
\[
i_*(x)i_*(y)=i_*(x\,i^*i_*(y)),
\]
结合 4.1 即得结论.

\begin{statement}{推论 4.6}
设 $S$ 为概形, $E$ 为秩恒为 $r$ 的局部自由 $\mathcal O_S$-模, $g:X=V(\check E)\to S$ 为截面层是 $E$ 的向量丛. 对 $X$ 在 $S$ 上的每个截面 $f$, 有
\[
f_*(1_S)=g^*(c_r(E))
\]
等式在 $H^{2r}(X,\mu_X^{\otimes r})$ 中成立.
\end{statement}

\emph{证明.} 由于 $g^*$ 是同构 (1.2), 且 $g\circ f=\mathrm{id}$, $f^*$ 是双射. 因此只需验证
\[
f^*f_*(1_S)=f^*g^*(c_r(E)).
\]
若以 $N$ 表示 $f$ 的法层, 则由 4.1, 最后这个等式等价于
\[
c_r(\check N)=c_r(E).
\]
事实上, 甚至有同构 $\check N\simeq E$, 因为将 (EGA IV 17.2.5) 的正合列应用于情形
\[
S\xrightarrow{f} X\xrightarrow{g} S,
\]
给出同构 $N\simeq f^*\Omega^1_{X/S}$, 且 $\Omega^1_{X/S}\simeq g^*\check E$ (EGA IV 16.4.8).

\begin{statement}{推论 4.7}
在 4.6 的假设下, 再假设 $S$ 在可分闭域上光滑且拟射影. 用 $\sigma_X$ 表示 $X$ 的零截面代数闭链, 令 $Y=f^{-1}(\sigma_X)$ 为 $f$ 的零点代数闭链. 则
\[
\cl_S(Y)=c_r(E)
\]
等式在 $H^{2r}(S,\mu_S^{\otimes r})$ 中成立.
\end{statement}

\emph{证明.} 由 IV,
\[
\cl_S(Y)=f^*\cl_X(\sigma_X).
\]
设 $s_0$ 为 $X$ 的零截面态射. 则
\[
\cl_X(\sigma_X)\stackrel{\mathrm{dfn}}{=}(s_0)_*(1_S)\stackrel{(4.6)}{=}g^*(c_r(E)),
\]
对首尾两项施以 $f^*$ 即得断言.

\begin{statement}{注 4.8}
设 $\ell$ 为与 $S$ 的剩余特征不同的素数. 直接取投射极限, 便由 4.6 和 4.7 得到 $\ell$ 进上同调中的相应断言.
\end{statement}

现在设 $\ell$ 为素数, $X$ 为可分闭域 $k$ 上的固有概形, $k$ 的特征与 $\ell$ 互素. 回顾 (IV), $X$ 的 $\ell$ 进 Euler--Poincaré 示性数是整数
\[
\operatorname{EP}_\ell(X)=\sum_i(-1)^i\dim H^i(X,\mathbb Q_\ell).
\]

\begin{statement}{推论 4.9}
设 $X$ 为可分闭域 $k$ 上的光滑固有连通概形, 维数为 $n$, 且 $\operatorname{char}(k)\ne\ell$. 若以 $p:X\to k$ 表示典范投影, 则
\[
\operatorname{EP}_\ell(X)=p_*c_n(\check\Omega^1_X)
\]
等式在 $\mathbb Q_\ell=H^0(k,\mathbb Q_\ell)$ 中成立.
\end{statement}

\emph{证明.} 设 $j:X\to X\times X$ 为 $X$ 的对角浸入, $q:X\times X\to k$ 为典范投影. 对态射 $j$ 应用 4.1 得到
\[
c_n(\check\Omega^1_X)=j^*j_*(1_X),
\]
所以
\[
p_*(c_n(\check\Omega^1_X))=p_*j^*j_*(1_X)=q_*j_*j^*j_*(1_X).
\]
利用 $j$ 的投影公式, 得到
\[
p_*(c_n(\check\Omega^1_X))=q_*(j_*(1_X)\cup j_*(1_X)),
\]
其中由 IV, 右边等于 $X$ 的 $\ell$ 进 Euler--Poincaré 示性数.

\subsubsection*{4.10}
当 $X$ 在 $k$ 上射影时, 上面引自 IV 的结果表明, $\operatorname{EP}_\ell(X)$ 在 $X\times X$ 的 Chow 环中等于对角代数闭链的自交数, 因而与素数 $\ell\ne\operatorname{char}(k)$ 无关. 所以可以不指明整数 $\ell$, 直接称为 $X$ 的 Euler--Poincaré 示性数 $\operatorname{EP}(X)$.

\begin{statement}{命题 4.11}
设 $X$ 为可分闭域 $k$ 上的 $n$ 维光滑射影连通概形. 则
\[
\operatorname{EP}(X)=\sum_{p,q}(-1)^{p+q}\dim H^p(X,\Omega_X^q).
\]
\end{statement}

\emph{证明.} 使用 4.9 的证明中的记号, 并令 $J$ 为 $X$ 的对角浸入的理想. 如 4.10 所述, 在 $X\times X$ 的 Chow 环中, $\operatorname{EP}(X)$ 是对角代数闭链 $\Delta$ 与自身的交积 $\Delta\cdot\Delta$. 为计算它, 可以忽略挠元, 因而在 $\operatorname{Gr}_{\mathrm{top}}^\bullet(X)$ 中计算. 这说明 $\operatorname{EP}(X)$ 是 $\mathcal O_{X\times X}/J$ 的类的平方在正像态射
\[
q_*:K(X\times X)\to K(k)=\mathbb Z
\]
下的像. 而根据 (SGA 6 VII 2.5), 这个平方等于
\[
\sum_i(-1)^i\lambda^i([\Omega_X^1]),
\]
所需公式立即由 $q_*$ 的定义得到.

\section*{5. 旗概形}

设 $S$ 为概形, $E$ 为秩 $r$ 的局部自由 $\mathcal O_S$-模, $(p_i)_{1\le i\le m}$ 为 $m$ 个正整数, 满足
\[
p_1+p_2+\cdots+p_m=r.
\]
回顾 $E$ 的 $(p_1,\ldots,p_m)$ 型旗是 $E$ 的有限滤过
\[
0=E_m\subset E_{m-1}\subset\cdots\subset E_0=E
\]
使得对每个整数 $i\in[1,m]$, 商 $E_{i-1}/E_i$ 局部自由且秩为 $p_i$.

定义函子 $(\operatorname{Sch}/S)^\circ\to(\operatorname{Ens})$: 它给每个 $S$-概形 $u:T\to S$ 配以 $u^*(E)$ 的 $(p_1,\ldots,p_m)$ 型旗的集合, 给每个 $S$-概形态射配以自然意义下旗的逆像. 已证明 (EGA I, 第 2 版) 此函子可表, 任何代表它的对象称为 $E$ 的 $(p_1,\ldots,p_m)$ 型旗的 $S$-概形.

以下用 $f:X\to S$ 表示 $E$ 的 $(p_1,\ldots,p_m)$ 型旗的 $S$-概形. 此外, 用 $D$ 表示 $E$ 的 $(1,1,\ldots,1)$ 型旗概形; 不会混淆时, 直接称为 $E$ 的旗概形.

由定义显然 $\mathcal O_X$-模 $f^*(E)$ 具有典范有限滤过
\[
0=F_m\subset F_{m-1}\subset\cdots\subset F_0=f^*(E),
\]
其中商 $E_i=F_{i-1}/F_i$ $(1\le i\le m)$ 局部自由, 秩为 $p_i$. 对每个整数 $i\in[1,m]$, 设
\[
g_i:D_i\to X
\]
为 $E_i$ 的旗的 $X$-概形, $g:D_1\times_X\cdots\times_X D_m\to X$ 为各 $X$-概形 $D_i$ 的纤维积. 显然复合
\[
h=f\circ g:\prod_iD_i\to S
\tag{5.1}
\]
与 $E$ 的旗的 $S$-概形在 $S$ 上同构; 因此在相差 $S$-同构的意义下可以写成
\[
D=\prod_iD_i.
\]

\begin{statement}{命题 5.2}
在 $D^b(S,A)$ 中有同构
\[
\mathrm Rf_*(A)\simeq\bigoplus_i R^{2i}f_*(A)[-2i].
\]
\end{statement}

\emph{证明.} 由于 $h$ 是有限多个形如
\[
P_T(F)\to T\qquad(F\text{ 为局部自由 }\mathcal O_T\text{-模}),
\tag{Q}
\]
的态射的复合, 由 (2.2.1) 得到
\[
\mathrm Rh_*(A)\simeq\bigoplus_j R^jh_*(A)[-j].
\]
同理, 还有
\[
\mathrm Rg_*(A)\simeq\bigoplus_j R^jg_*(A)[-j].
\]
利用 ([6] 2.16), 得到同构
\[
\mathrm Rf_*(g_*A)\simeq\bigoplus_j R^j f_*(g_*A)[-j].
\]
由于 $g$ 光滑且几何纤维连通, $A_X\simeq g_*(A_D)$; 因而只需证明
\[
R^{2i+1}f_*(A)=0\qquad (i\in\mathbb N).
\tag{5.2.1}
\]
为此, 利用纤维函子族的保守性 (SGA 4 VIII 3.5), 立即归约到 $S$ 是可分闭域 $k$ 的谱的情形. 此时, 由于 $g$ 为射影丛的复合, (2.2.4) 表明逆像映射
\[
H^{2i+1}(X,A_X)\to H^{2i+1}(D,A_D)
\]
为单态射, 故可假设 $X=D$. 在此情形, 投影 $D\to k$ 本身也是射影丛的复合, 所以由 (2.2.4) 归纳即得断言.

\begin{statement}{推论 5.3}
对每个概形态射 $u:S\to T$, Leray 谱序列
\[
E_2^{pq}=R^p u_*\,R^q f_*(A)\Rightarrow R^{p+q}(u\circ f)_*(A)
\]
退化.
\end{statement}

\emph{证明.} 参见 [6] 1.2.

现在明确 $X$ 的上同调的乘法结构, 先考察特殊情形. 用前述记号, 得到环态射
\[
\alpha:A^*(S)[T_{ij_i}]_{1\le i\le m,\,1\le j_i\le p_i}\to A^*(X)
\]
它通过逆像 $f^*$ 将 $A^*(S)$ 映入 $A^*(X)$, 对上述每对 $(i,j_i)$, 将元素 $T_{ij_i}$ 送到第 $j_i$ 个 Chern 类 $c_{j_i}(E_i)$. 根据 (3.4 (iii)), 对每个整数 $p$, 有
\[
c_p(f^*E)=\sum_{i_1+i_2+\cdots+i_m=p}c_{i_1}(E_1)\cdots c_{i_m}(E_m),
\]
因而由 Chern 类的函子性, 此态射在以各个 $p$ 所对应的元素
\[
c_p(E)-\sum_{i_1+\cdots+i_m=p}T_{1,i_1}\cdots T_{m,i_m}
\]
为生成元的理想 $J_X$ 上为零. 因此有环态射
\[
\varphi_X:\bigl(A^*(S)[T_{ij_i}]_{(i,j_i)}\bigr)/J_X\to A^*(X).
\]

\begin{statement}{命题 5.4}
上述环态射 $\varphi_X$ 为同构. $A^*(S)$-模 $A^*(X)$ 是秩为下式的自由模:
\[
\frac{r!}{p_1!p_2!\cdots p_m!}.
\]
\end{statement}

\emph{证明.} 遵循 Grothendieck 的报告 (Séminaire Chevalley 1958, Anneaux de Chow et applications, Exp. IV). 将用到下列引理 (参见该文第 19 页引理 1).

\begin{statement}{引理 5.4.1}
a) 设 $R$ 为交换环, $(a_i)_{1\le i\le r}$ 为 $R$ 的 $r$ 个元素. 设 $\sigma_i$ 为元素 $(U_i)_{1\le i\le r}$ 的第 $i$ 个初等对称函数, $I$ 为多项式环 $R[U_1,U_2,\ldots,U_r]$ 中由元素 $\sigma_i-a_i$ 生成的理想. 则 $R$-代数
\[
R[U_1,U_2,\ldots,U_r]/I
\]
为秩 $r!$ 的自由 $R$-模.

b) 设 $r,p_1,p_2,\ldots,p_m$ 为满足 $p_1+\cdots+p_m=r$ 的正整数. 则多项式环态射
\[
V=R[T_{i,j_i}]_{1\le i\le m,\,1\le j_i\le p_i}
\longrightarrow R[U_i]_{1\le i\le r}=W
\]
是单射并使 $W$ 成为秩 $p_1!\cdots p_m!$ 的自由 $V$-模; 此态射将 $T_{i,j_i}$ 送到元素
\[
(U_{p_1+\cdots+p_{i-1}+s})_{1\le s\le p_i}
\]
的第 $j_i$ 个初等对称函数.

特别地, $R$-代数态射
\[
R[T_1,\ldots,T_r]\to R[U_1,\ldots,U_r]
\]
将 $T_i$ 送到各 $U_j$ 的第 $i$ 个初等对称函数; 它是单射, 并使右边成为左边上的秩 $r!$ 的自由模.
\end{statement}

说明如何从 (5.4.1) 得到 (5.4). 首先假设 $X$ 为 $E$ 的旗概形 $D$, 为简便记 $T_{i1}=U_i$. 由于典范投影 $D\to S$ 是 $r-1$ 个形如
\[
P_{Y_i}(M_i)\to Y_i\quad(1\le i\le r-1),
\]
的态射的复合, 其中 $M_i$ 为秩 $r-i+1$ 的局部自由 $\mathcal O_{Y_i}$-模, 由 (2.2.6), $A^*(D)$ 是秩 $r!$ 的自由 $A^*(S)$-模, 且态射 $\varphi_D$ 为满射. 而引理 5.4.1 a) 表明 $\varphi_D$ 的源也是秩 $r!$ 的自由 $A^*(S)$-模, 因而断言在 $X=D$ 时成立.

一般情形下, 显然下图交换:
\[
\begin{tikzcd}[column sep=large,row sep=large]
A^*(X) \arrow[r,"g^*"] & A^*(D)\\
C=A^*(S)[T_{i,j_i}]/J_X \arrow[u,"\varphi_X"] \arrow[r,"\gamma"'] &
A^*(S)[U_k]/J_D \arrow[u,"\varphi_D"']
\end{tikzcd}
\]
这里映射 $\gamma$ 是将下述 $A^*(S)$-代数态射取商得到的:
\[
A^*(S)[T_{i,j_i}]\to A^*(S)[U_k]
\]
它将 $T_{i,j_i}$ 送到元素
\[
(U_{p_1+\cdots+p_{i-1}+s})_{1\le s\le p_i}.
\]
的第 $j_i$ 个初等对称函数. 由于 $\varphi_D$ 是同构, 5.4.1 b) 表明 $A^*(D)$ 是秩 $p_1!\cdots p_m!$ 的自由 $C$-模. 另一方面, 将 $g$ 写成旗的 $X$-概形的纤维积, 可见 $A^*(D)$ 也是秩 $p_1!\cdots p_m!$ 的自由 $A^*(X)$-模. 立即得到 $\varphi_X$ 是同构.

为陈述下一命题, 约定对每个概形 $Y$ 及每个 $A_Y$-模 $M$, 令
\[
Q(Y,M)=\bigoplus_{n,i}H^n(Y,M\otimes\mu_Y^{\otimes i}).
\]
$A$-模 $Q(Y,M)$ 通过杯积典范地带有 $A^*(Y)$-模结构. 若 $u:Z\to Y$ 为概形态射, 逆像态射
\[
u^*:Q(Y,M)\to Q(Z,u^*M)
\]
与 $Q(Y,M)$ 的 $A^*(Y)$-模结构和 $Q(Z,u^*M)$ 的 $A^*(Z)$-模结构相容. 因而导出称为典范态射的 $A^*(Z)$-模态射
\[
A^*(Z)\otimes_{A^*(Y)}Q(Y,M)\to Q(Z,u^*M).
\]

\begin{statement}{命题 5.5}
设 $S$ 为概形, $E$ 为秩 $r$ 的局部自由 $\mathcal O_S$-模. 设 $f:X\to S$ 为 $E$ 的 $(p_1,\ldots,p_m)$ 型旗的 $S$-概形. 则对每个 $A_S$-模 $L$, 典范态射
\[
\alpha:A^*(X)\otimes_{A^*(S)}Q(S,L)\to Q(X,f^*L)
\]
是同构.
\end{statement}

\emph{证明.} 当 $X=P_S(E)$ 时, 这是 (2.2.4) 和 (2.2.6) 的直接推论. 从而 $X=D$ 时也成立. 一般情形下, 由于典范态射 $g:D\to X$ 和 $h:D\to S$ 为射影丛的复合, 典范态射
\[
A^*(D)\otimes_{A^*(S)}Q(S,L)\to Q(D,h^*L)
\]
与
\[
A^*(D)\otimes_{A^*(X)}Q(X,g^*L)\to Q(D,g^*L)
\]
是同构. 所以 $\mathrm{id}_{A^*(D)}\otimes\alpha$ 为同构, 又因 $A^*(D)$ 在 $A^*(X)$ 上自由, 得到结论.

现在给出 (5.4) 与 (5.5) 的相对变体. 用前述记号, 设 $u:S\to T$ 为概形态射. 令
\[
A^*(X/T)=\bigoplus_i R^{2i}(u\circ f)_*(A_X)
\]
类似地,
\[
A^*(S/T)=\bigoplus_i R^{2i}u_*(A_S).
\]
最后, 对每个 $A_S$-模 $L$, 令
\[
Q(S/T,L)=\bigoplus_{n,i}R^n u_*(L\otimes\mu_S^{\otimes i})
\]
以及
\[
Q(X/T,f^*L)=\bigoplus_{n,i}R^n(u\circ f)_*(f^*L\otimes\mu_X^{\otimes i}).
\]

对每对满足 $1\le i\le m$, $1\le j\le p_i$ 的整数 $(i,j)$, Chern 类 $c_j(E_i)$ 在典范态射
\[
H^{2j}(X,\mu_X^{\otimes j})
\to H^0(T,R^{2j}(u\circ f)_*(\mu_X^{\otimes j}))
\]
下的像给出 $A_T$-态射
\[
c_{ij}:A_T\to R^{2j}(u\circ f)_*(\mu_X^{\otimes j}).
\]
各 $c_{ij}$ 及逆像 $f^*$ 确定 $A^*(S/T)$-代数态射
\[
\theta_1:A^*(S/T)[T_{i,j_i}]_{(i,j_i)}\to A^*(X/T).
\]
用 $c_i$ 表示 Chern 类 $c_i(E)$ 确定的 $A^*(S/T)$ 的截面; 如同 (5.4) 中那样, 可见 $\theta_1$ 在由下列关系局部生成的理想 $J$ 上为零:
\[
\sum_{i_1+\cdots+i_m=p}T_{1,i_1}\cdots T_{m,i_m}=c_p\qquad(p\in\mathbb N).
\]
因此有 $A^*(S/T)$-代数态射
\[
\theta:A^*(S/T)[T_{i,j_i}]_{(i,j_i)}/J\to A^*(X/T).
\]

\begin{statement}{命题 5.6}
沿用上述记号:
\begin{enumerate}[label=\alph*)]
\item $A^*(S/T)$-代数态射
\[
\theta:A^*(S/T)[T_{i,j_i}]_{(i,j_i)}/J\to A^*(X/T)
\]
是同构.
\item 对每个 $A_S$-模 $L$, 典范 $A^*(X/T)$-模态射
\[
A^*(X/T)\otimes_{A^*(S/T)}Q(S/T,L)\to Q(X/T,f^*L)
\]
是同构.
\end{enumerate}
\end{statement}

\emph{证明.} 分别对 (5.4) 和 (5.5) 作层化即可.

\section*{6. 群概形}

先给出由挠子扭曲对象的一个一般性质.

\begin{statement}{命题 6.1}
设 $S$ 为概形, $G$ 为拟紧光滑且纤维连通的 $S$-群概形, $X$ 为带有 $G$ 的左作用的 $S$-概形. 若 $f:E\to S$ 为 $G$ 下的右 $S$-挠子, 用 $X_E$ 表示借助 $E$ 和 $G$ 在 $X$ 上的作用, 从 $X$ 扭曲得到的 $S$-概形. 设 $p:X\to S$, $q:X_E\to S$ 为典范投影. 在这些条件下, 对每个 $A_S$-模 $M$ 及每个整数 $i$, $A_S$-模 $R^iq_*(q^*M)$ 典范同构于 $R^ip_*(p^*M)$.
\end{statement}

\emph{证明.}
\[
\begin{tikzcd}[column sep=large,row sep=large]
E\times_S X \arrow[r,"\varphi"] \arrow[d,"\pr_1"'] & X_E \arrow[d,"q"]\\
E \arrow[r,"f"'] & S
\end{tikzcd}
\qquad
\begin{tikzcd}[column sep=large,row sep=large]
E\times_S X \arrow[r,"\pr_2"] \arrow[d,"\pr_1"'] & X \arrow[d,"p"]\\
E \arrow[r,"f"'] & S .
\end{tikzcd}
\]
已知图 $(D_1)$ 为笛卡尔图, 其中 $\varphi$ 表示典范取商映射. 因为 $f$ 光滑, 光滑基变换定理给出同构
\[
u:f^*R^iq_*(q^*M)\xrightarrow{\sim}R^i(\pr_1)_*(\varphi^*q^*M).
\]
类似地, 笛卡尔图 $(D_2)$ 给出同构
\[
v:f^*R^ip_*(p^*M)\xrightarrow{\sim}R^i(\pr_1)_*((\pr_2)^*p^*M).
\]
而 $p\circ\pr_2=q\circ\varphi$, 因此将 $\varphi^*\circ q^*$ 与 $(\pr_2)^*\circ p^*$ 自然等同后, 得到同构
\[
v^{-1}\circ u:f^*R^iq_*(q^*M)\xrightarrow{\sim} f^*R^ip_*(p^*M).
\]
由于 $f$ 的几何纤维非空且连通, 根据 (SGA 4 XII 6.5 (i)), 伴随态射
\[
\operatorname{adj}:\mathrm{id}\to f_*f^*
\]
是同构, 所以
\[
(\operatorname{adj})^{-1}\circ f_*(v^{-1}\circ u)\circ(\operatorname{adj})
\]
是所需同构.

\begin{statement}{命题 6.2}
(Mme Raynaud [5].) 设 $S$ 为概形, $p:G\to S$ 为约化 $S$-群概形 (SGA 3 XIX). 对每个局部常值可构造 $A_S$-模 $M$ 及每个整数 $i$, $A_S$-模 $R^ip_*(p^*M)$ 局部常值且可构造. 此外, 对每个确定笛卡尔图
\[
\begin{tikzcd}[column sep=large,row sep=large]
G' \arrow[r,"u'"] \arrow[d,"p'"'] & G \arrow[d,"p"]\\
S' \arrow[r,"u"'] & S,
\end{tikzcd}
\]
的态射 $u:S'\to S$, 基变换态射
\[
u^*R^ip_*(p^*M)\to R^ip'_*(p\circ u')^*(M)
\]
是同构.
\end{statement}

\emph{证明.} 命题的两个断言对 $S$ 的平展拓扑都是局部的. 由 (SGA 3 XII 2.3), 可假设 $G$ 分裂. 设 $B$ 为 $G$ 的一个 (SGA 3 XXII 5.1.1) 所指定类型的 Borel 子群, 则由 (SGA 3 XXII 5.8.5), 商 $G/B$ 存在. 设 $q:G\to G/B$, $r:G/B\to S$ 为典范态射. 由于 $r$ 固有且光滑, 它满足固有基变换定理, 且函子 $R^ir_*$ 将局部常值可构造阿贝尔层送到同类层 (SGA 4 XVI 2.2). 此外, 根据 (SGA 4 XXII 5.5.1), $S$-概形 $B$ 同构于有限个向量丛与去掉零截面的向量丛的乘积. 因而立即由 (1.1) 和 (1.2) 得到: 若 $t:B\to S$ 表示典范态射, $M$ 为局部常值可构造 $A_S$-模, 则层 $R^it_*(t^*M)$ 局部常值、可构造, 且在基变换下稳定. 现在设 $M$ 为局部常值可构造 $A_S$-模. 上述讨论和 (6.1) 表明层 $R^iq_*(p^*M)$ 局部常值且可构造; 利用 Leray 谱序列
\[
R^ir_*(R^jq_*)(p^*M)\Rightarrow R^{i+j}p_*(p^*M),
\]
便对层 $R^ip_*(p^*M)$ 得到相同结论. 最后, 由 (6.1), 层 $R^iq_*(p^*M)$ 在基变换下稳定; 对 $r$ 使用固有基变换定理, 即得层 $R^ip_*(p^*M)$ 也在基变换下稳定.

\begin{statement}{注 6.3}
由于每个约化 $S$-群在平展拓扑下局部都是某个约化 $\operatorname{Spec}(\mathbb Z)$-概形沿典范态射 $S\to\operatorname{Spec}(\mathbb Z)$ 的逆像 (SGA 3 XXIII 5.7 和 XXV 1.3), (6.2) 表明: 在代数闭域上计算约化群概形的与特征互素的挠系数上同调, 可归约到在复数域上计算同类型约化群概形的上同调这一超越问题. 由比较定理 (SGA 4 XI 4.4), 后者可用代数拓扑的方法求得.
\end{statement}

\begin{statement}{推论 6.4}
设 $S$ 为概形, $p:G\to S$ 为 $S$-群概形, 它是有限平展 $S$-群概形 $F$ 以约化 $S$-群概形 $H$ 为核的扩张. 则 (6.2) 的结论对每个局部常值可构造 $A_S$-模 $M$ 都成立.
\end{statement}

\emph{证明.} 作基变换 $F\to S$ 后, 可假设 $F$ 形如 $\Gamma_S$, 其中 $\Gamma$ 为通常的有限群. 此时有 $S$-概形的同构
\[
G\simeq\coprod_{\gamma\in\Gamma}H,
\]
结合 (6.2) 即得结论.

\begin{statement}{推论 6.5}
设 $S$ 为概形, $G$ 为约化 $S$-群概形, $H$ 为 $G$ 的约化 $S$-子群概形, 假设忠实平坦拟紧拓扑中 $G$ 关于 $H$ 的商层可表. 则若
\[
p:G/H\to S
\]
为典范态射, (6.2) 的结论对每个局部常值可构造 $A_S$-模 $M$ 均成立.
\end{statement}

\emph{证明.} 设 $f:G\to G/H$ 为典范态射; 它光滑, 因而是平展层的下降态射 (SGA 4 VIII 9.1). 对每个整数 $r\ge0$, 令
\[
G_r=G\times_{G/H}G\times_{G/H}\cdots\times_{G/H}G
\quad(r+1\text{ 次}),
\]
由 SGA 3 V 1.5, 它典范地等同于
\[
G\times_S H\times_S\cdots\times_S H,
\]
其中 $H$ 出现 $r$ 次. 设 $f_r:G_r\to G/H$, $g_r:G_r\to S$ 为典范态射, 特别有 $f_0=f$. 最后, 对每个整数 $q\ge0$, 用 $H^q(M)$ 表示复形
\[
R^q(g_0)_*(f_0^*M)\to R^q(g_1)_*(f_1^*M)\rightrightarrows\cdots\to R^q(g_r)_*(f_r^*M)\cdots .
\]
% SGA-SRC-000365: the scan leaves the citation blank; reconstructed from the mathematics.
根据 (SGA 4, 第 V bis 报告, 2.5.5 和 2.5.6), 有下降谱序列, 又称 Čech 谱序列,
\[
E(M):E_2^{pq}=H^p(H^q(M))\Rightarrow E^n=R^np_*(M).
\]
由于各概形 $G_r$ 都是约化群概形, 由 6.2, 各项 $E_2^{pq}$ 局部常值、可构造, 且与任意基变换相容. 立即得到各 $E^n$ 也如此.

\section*{7. 完全交}

本节给定可分闭域 $k$. 若 $\ell$ 为不同于 $\operatorname{char}(k)$ 的素数, 则对每个固有 $k$-概形 $X$ 和每个整数 $i\ge0$, 用
\[
b_\ell^i(X)\quad\text{或}\quad b^i(X)
\]
表示 $X$ 的 $\mathbb Z_\ell$ 系数第 $i$ 个 Betti 数; 后一种记号用于不会混淆时. 若 $X$ 在 $k$ 上射影, 则 $X$ 的超平面截面指向量丛 $V(\check{\mathcal L})$ 的一个与零截面横截的截面的零点集, 其中 $\mathcal L$ 为丰沛可逆 $\mathcal O_X$-模.

\begin{statement}{定理 7.1}
(Lefschetz 定理.) 设 $X$ 为 $n$ 维光滑连通射影 $k$-概形, $Y$ 为 $X$ 的 $d$ 个超平面截面的交所成的子概形, $L$ 为局部常值可构造 $A_X$-模. 则:
\begin{enumerate}[label=(\roman*)]
\item 逆像同态
\[
H^i(X,L)\to H^i(Y,L|Y)
\]
当 $i\le n-d-1$ 时是同构, 当 $i=n-d$ 时是单态射.
\item 假设 $Y$ 在 $k$ 上光滑, 维数为 $n-d$. 则 Gysin 同态
\[
H^i(Y,(L|Y)\otimes\mu_Y^{\otimes -d})\to H^{i+2d}(X,L)
\]
当 $i\ge n-d+1$ 时是同构, 当 $i=n-d$ 时是满态射.
\end{enumerate}
\end{statement}

\emph{证明.} 众所周知, 超平面截面的补集为仿射开集, 因而 $U=X-Y$ 是 $d$ 个仿射开集的并. 断言 (i) 立即由上同调正合列
\[
\cdots\to H_!^i(U,L|U)\to H^i(X,L)\to H^i(Y,L|Y)
\xrightarrow{\partial} H_!^{i+1}(U,L|U)\to\cdots
\]
及以下引理得到.

\begin{statement}{引理 7.1.1}
设 $U$ 为维数 $n$ 的有限型 $k$-概形, 它是 $d$ 个仿射开子集的并, $F$ 为 $A_U$-模. 则:
\begin{enumerate}[label=\alph*)]
\item 若 $F$ 可构造,
\[
H^i(U,F)=0\qquad\text{对 }i\ge n+d;
\]
\item 若 $F$ 局部常值, 且 $U$ 在 $k$ 上光滑,
\[
H^i_c(U,F)=0\qquad\text{对 }i\le n-d.
\]
\end{enumerate}
\end{statement}

当 $d=1$ 时, 断言 a) 就是 (SGA 4 XIV 3.2). 一般情形下, 以 $\mathfrak U$ 表示由 $d$ 个仿射开子集组成的 $U$ 的覆盖, 从 Čech 谱序列
\[
H^p(\mathfrak U,\mathcal H^q(F))\Rightarrow H^{p+q}(U,F).
\]
得到 a). 下面证明 b). 由于函子 $H_c^i(U,\cdot)$ 与滤过归纳极限交换 (SGA 4 VII 3.3), 可假设 $F$ 可构造; 此时由 a) 及 Poincaré 对偶 (SGA 4 XVIII) 即得断言.

最后, 利用 Gysin 态射通过典范同构
\[
H^{i-2d}(Y,(L|Y)\otimes\mu_Y^{\otimes -d})\simeq H^i_Y(X,L),
\tag{7.1.2}
\]
给出的定义, 断言 (ii) 由 7.1.1 a) 以及开子集及其补集的上同调正合列
\[
\cdots\to H^{i-1}(U,L|U)\xrightarrow{\partial}H^i_Y(X,L)\to H^i(X,L)\to H^i(U,L|U)\to\cdots .
\]

\begin{statement}{注 7.1.3}
事实上, 由切除性质, 当 $X$ 光滑、$L$ 可构造且局部常值的假设仅在 $Y$ 的某个邻域内成立时, 典范同构 (7.1.2) 仍可推广; 从而 7.1 (ii) 也推广到这一背景.

实质上更一般的表述参见 (SGA 2 XIV 4.6), 但那里的证明只借助奇点消解, 因此暂时限于特征 $0$.
\end{statement}

\begin{statement}{推论 7.2}
设 $X$ 为 $n$ 维光滑连通射影 $k$-概形, $Y$ 为 $X$ 的 $d$ 个超平面截面的交所成的子概形, $L$ 为可构造扭常值 $\mathbb Z_\ell$-层 $(\ell\ne\operatorname{car}(k))$. 则:
\begin{enumerate}[label=(\roman*)]
\item 逆像同态
\[
\alpha_i:H^i(X,L)\to H^i(Y,L|Y)
\]
当 $i\le n-d-1$ 时为同构, 当 $i=n-d$ 时为单态射. 此外, 若 $L$ 可构造且局部自由, 则 $\alpha_{n-d}$ 的余核无挠;
\item 假设 $Y$ 在 $k$ 上光滑, 维数为 $n-d$. 则 Gysin 同态
\[
\beta_i:H^i(Y,(L|Y)(-d))\to H^{i+2d}(X,L)
\]
当 $i\ge n-d+1$ 时为同构, 当 $i=n-d$ 时为满态射. 此外, 若 $L$ 可构造且局部自由, 则满态射 $\beta_{n-d}$ 分裂, 且其核无挠.
\end{enumerate}
\end{statement}

\emph{证明.} 证明 (i), 断言 (ii) 随后由 Poincaré 对偶得到. 一般情形下, 从 7.1 (i) 直接取极限便得到结论. 当 $L$ 为局部自由 $\mathbb Z_\ell$-层时, 对 $D(\mathbb Z_\ell)$ 中的逆像态射
\[
R\Gamma(X,L)\to R\Gamma(Y,L|Y)
\]
应用以下引理; 假设 (iii) 由 7.1 (i) 满足.

\begin{statement}{引理 7.2.1}
设 $V$ 为离散赋值环, $k_V$ 为其剩余域, $u:K\to L$ 为 $D(V)$ 中的态射. 假设 $K,L$ 的上同调为有限型, 用 $M$ 表示 $u$ 的映射柱. 则对每个 $m\in\mathbb Z$, 下列条件等价:
\begin{enumerate}[label=(\roman*)]
\item 当 $p\le m-1$ 时有 $H^p(M)=0$, 且 $H^m(M)$ 无挠;
\item 对每个 $V$-模 $N$, 态射
\[
H^p(u\otimes^L\operatorname{id}_N):H^p(K\otimes^L N)\to H^p(L\otimes^L N)
\]
当 $p\le m-1$ 时是同构, 当 $p=m$ 时是单态射;
\item[(ii bis)] 对每个 $V$-模 $N$, 有
\[
H^p(M\otimes^L N)=0\qquad(p\le m-1);
\]
\item 态射
\[
H^p(u\otimes^L\operatorname{id}):H^p(K\otimes^L k_V)\to H^p(L\otimes^L k_V)
\]
当 $p\le m-1$ 时是同构, 当 $p=m$ 时是单态射;
\item[(iii bis)] 有
\[
H^p(M\otimes^L k_V)=0\qquad(p\le m-1).
\]
\end{enumerate}
此外, 这些等价条件蕴含:
\begin{enumerate}[label=(\roman* bis)]
\item 态射
\[
H^p(u):H^p(K)\to H^p(L)
\]
当 $p\le m-1$ 时是同构, 当 $p=m$ 时是具有自由余核的单态射.
\end{enumerate}
\end{statement}

\emph{证明.} 下列蕴含关系图是显然的:
\[
\begin{array}{ccc}
\text{(ii)}&\Leftrightarrow&\text{(ii bis)}\\
\Downarrow&&\\
\text{(iii)}&\Leftrightarrow&\text{(iii bis)}
\end{array}
\qquad\qquad
\text{(i)}\Rightarrow\text{(i bis).}
\]
将分别证明 (i)$\Rightarrow$(ii bis) 及 (iii bis)$\Rightarrow$(i), 这就足够了. 蕴含 (i)$\Rightarrow$(ii bis) 直接由 Künneth 正合列
\[
0\to H^p(M)\otimes_V N\to H^p(M\otimes^L N)\to \operatorname{Tor}_1^V(H^{p+1}(M),N)\to0.
\]
得到. 对 $N=k_V$ 使用同一正合列, 可见若 (iii bis) 成立, 则
\[
H^p(M)\otimes k_V=0\quad(p\le m-1),\qquad
\operatorname{Tor}_1^V(H^m(M),k_V)=0.
\]
由 Nakayama 引理及平坦性局部判别准则 (EGA $0_{\mathrm{III}}$ 10.1.3) 得到 (i).

在上述假设下, 可见知道 $X$ 的 $\mathbb Q_\ell$ 系数上同调群, 就能知道 $Y$ 的上同调群, 唯一例外是 $n-1$ 维. 在这些条件下, 余下那个群的计算立即归约为 $Y$ 的 Euler--Poincaré 示性数的计算. 下面在稍一般的背景下计算它.

\begin{statement}{命题 7.3}
设 $X$ 为 $n$ 维固有光滑连通 $k$-概形, $E$ 为秩恒为 $r$ 的局部自由 $\mathcal O_X$-模, $Y$ 为 $V(\check E)$ 的一个与零截面横截的截面的零点集. 设 $p:X\to k$ 为典范态射, $p_*$ 为相应的 Gysin 态射. 则
\[
\operatorname{EP}(Y)=p_*\bigl(c(\check\Omega_X^1)c(E)^{-1}c_r(E)\bigr)
\]
等式在 $H^0(k,\mathbb Q_\ell)=\mathbb Q_\ell$ 中成立.
\end{statement}

\emph{证明.} 由 (EGA IV 17.13.2), $Y$ 光滑且纯维数为 $n-r$. 若 $i:Y\to X$ 表示典范浸入, 则 $i$ 的余法层是 $V(\check E)$ 的零截面的余法层 $\check E$ 沿 $i$ 的逆像. 因而有正合列
\[
0\to i^*(\check E)\to i^*(\Omega_X^1)\to\Omega_Y^1\to0,
\tag{7.3.1}
\]
由此立即得到
\[
c(\check\Omega_Y^1)=i^*\bigl(c(\check\Omega_X^1)c(E)^{-1}\bigr).
\tag{7.3.2}
\]
由 (3.10), 因此得到
\[
\operatorname{EP}(Y)=p_*i_*i^*\bigl(c(\check\Omega_X^1)c(E)^{-1}\bigr),
\]
或利用 $i$ 的投影公式, 得到
\[
\operatorname{EP}(Y)=p_*\bigl(c(\check\Omega_X^1)c(E)^{-1}i_*(1_Y)\bigr).
\]
命题的公式立即由 (4.7) 得到.

\begin{statement}{推论 7.4}
设 $X$ 为 $n$ 维固有光滑连通 $k$-概形, $Y$ 为 $X$ 的 $n-1$ 维光滑闭子概形, 因而由某个可逆 $\mathcal O_X$-模 $L$ 的一个截面定义. 设 $p$ 为 $X$ 到 $k$ 的典范态射, $p_*$ 为相应的 Gysin 态射. 则
\[
\operatorname{EP}(X)-\operatorname{EP}(Y)=p_*\bigl(c(\check\Omega_X^1)c(L)^{-1}\bigr)
\tag{7.4.1}
\]
等式在 $H^0(k,\mathbb Q_\ell)=\mathbb Q_\ell$ 中成立. 此外,
\[
b^{n-1}(Y)=2b^{\,n-1}(X)-b^n(X)+(-1)^n p_*\bigl(c(\check\Omega_X^1)c(L)^{-1}\bigr).
\tag{7.4.2}
\]
\end{statement}

\emph{证明.} 公式 (7.4.1) 由 7.3 及 4.2, 结合显然的恒等式
\[
1-c(L)^{-1}c_1(L)=c(L)^{-1}.
\]
得到. 利用 (7.2), 可见
\[
b^{n-1}(Y)=b^{n-1}(X)-b^n(X)+b^{n+1}(X)+(-1)^n(\operatorname{EP}(X)-\operatorname{EP}(Y)).
\]
而由 Poincaré 对偶, $b^{n-1}(X)=b^{n+1}(X)$, 从而得到 (7.4.2).

\begin{statement}{推论 7.5}
设 $X$ 为 $\mathbb P_k^{\,n+r}$ 的 $n$ 维光滑连通子概形, 它是 $\mathbb P_k^{\,n+r}$ 的 $r$ 个除子 $X_1,X_2,\ldots,X_r$ 的整体完全交. 对每个 $i\in\{1,\ldots,r\}$, 设 $d_i$ 为 $X_i$ 的次数. 则:
\begin{enumerate}[label=(\roman*)]
\item 上同调群 $H^p(X,\mathbb Z_\ell)$ 为有限型自由 $\mathbb Z_\ell$-模;
\item 若整数 $p\le 2n$ 且 $p\ne n$, 则
\[
H^p(X,\mathbb Z_\ell)=
\begin{cases}
0,&\text{若 }p\text{ 为奇数},\\
\mathbb Z_\ell,&\text{若 }p\text{ 为偶数};
\end{cases}
\]
\item $X$ 的 Euler--Poincaré 示性数为
\[
\operatorname{EP}(X)
=d_1\cdots d_r
\sum_{\substack{0\le i\le n\\ j_1+\cdots+j_r=n-i}}
(-1)^{\,n-i}\binom{n+r+1}{i}
d_1^{j_1}\cdots d_r^{j_r}.
\]
\end{enumerate}
\end{statement}

\emph{证明.} 断言 (i) 立即由 (7.2) 得到; 结合射影空间的上同调 (第 2 节), (ii) 也立即得到. 下面证明 (iii). 设 $P$ 为所在的射影空间, $p:P\to k$ 为典范投影, $p_*$ 为相应的 Gysin 态射, $\xi$ 为 $\mathcal O_P(1)$ 的类. 对每个 $i\in\{1,\ldots,r\}$, 概形 $X_i$ 由某个可逆 $\mathcal O_P$-模 $L_i$ 的截面定义, 所以 $X$ 是 $V(\check L_1\oplus\cdots\oplus\check L_r)$ 的一个与零截面横截的截面的零点集. 另一方面, 有正合列
\[
0\to\Omega_P^1\otimes\mathcal O_P(1)\to\mathcal O_P^{\,n+r+1}\to\mathcal O_P(1)\to0,
\]
所以
\[
c(\check\Omega_P^1)=(1+\xi)^{\,n+r+1}.
\]
公式 (7.3) 给出
\[
\operatorname{EP}(X)
=p_*\bigl(c(\check\Omega_P^1)c(\check L_1\oplus\cdots\oplus\check L_r)^{-1}
c_1(\check L_1)\cdots c_1(\check L_r)\bigr)
\]
\[
=p_*\bigl(d_1\cdots d_r\,\xi^r(1+\xi)^{\,n+r+1}
(1+d_1\xi)^{-1}\cdots(1+d_r\xi)^{-1}\bigr).
\]
利用 $p_*(\xi^{n+r})=1$ (2.2.2 c)), 不难得到命题中的公式.

\begin{statement}{注 7.6}
也可使用 (4.11), 转写 Hirzebruch 的计算 ([3], pp. 160--161), 得到 Euler--Poincaré 示性数的表达式; 这一计算使用 Riemann--Roch 定理. 关于更完整的计算和完全交的 Hodge 上同调的确定, 请读者参见 (SGA 7 XI).
\end{statement}

\subsection*{7.7}

类比超越情形, 在那里对常系数有另一个熟知的 Lefschetz 定理 ([4]), 自然猜想以下断言. 若 $X$ 为 $n$ 维光滑连通射影 $k$-概形, $\xi$ 为极丰沛可逆 $\mathcal O_X$-模在 $H^2(X,\mathbb Z_\ell(1))$ 中的类, 则对每个扭常值 $\mathbb Q_\ell$-层 $L$ 和每个整数 $p\le n$, 乘以 $\xi^{n-p}$ 的映射
\[
H^p(X,L)\to H^{2n-p}(X,L(n-p))
\]
为同构. 如 Deligne 所指出, 若不对 $L$ 加上附加假设, (7.7) 的断言不成立. 不过, Deligne 在 \emph{La conjecture de Weil II} 中证明: 当 $L$ 还满足纯性时, 这些断言成立.

这一猜想有下列推论, 同样在 $\operatorname{car}\ne0$ 时尚未证明.

设 $X$ 为 $n$ 维光滑连通射影 $k$-概形, $Y$ 为 $X$ 的光滑超平面截面, $L$ 为 $X$ 上的扭常值 $\mathbb Q_\ell$-层.
\begin{enumerate}[label=(\roman*)]
\item 逆像态射
\[
H^p(X,L)\to H^p(Y,L|Y)
\]
当 $p\ge n$ 时是满态射.
\item Gysin 态射
\[
H^p(Y,(L|Y)(-1))\to H^{p+2}(X,L)
\]
当 $p\le n-2$ 时是单态射.
\end{enumerate}

\emph{证明.} 设 $j:X\hookrightarrow\mathbb P^r$ 为闭浸入, $\xi_X$ (分别地 $\xi_Y$) 为 $\mathcal O_X(1)$ (分别地 $\mathcal O_Y(1)$) 在 $H^2(X,\mathbb Q_\ell(1))$ (分别地 $H^2(Y,\mathbb Q_\ell(1))$) 中的类. 用 $\alpha:Y\hookrightarrow X$ 表示典范嵌入. 对 $q\le n-1$, 下图
\[
\begin{array}{ccc}
H^q(X,L)&\xrightarrow{\alpha^*}&H^q(Y,L|Y)\\
\cup\,\xi_X^{\,n-1-q}\downarrow&&\downarrow\cup\,\xi_Y^{\,n-1-q}\\
H^{2n-2-q}(X,L(n-1-q))&\xrightarrow{\alpha^*}&
H^{2n-2-q}(Y,(L|Y)(n-1-q))
\end{array}
\]
交换, 而右列猜想为同构. 令 $p=2n-2-q$. 当 $q\le n-2$, 即 $p\ge n$ 时, 由 (7.2), 上行为同构, 从而得到 (i). 断言 (ii) 的证明类似, 这次考虑下图, 其中 $p\le n-1$:
\[
\begin{array}{ccc}
H^p(Y,(L|Y)(-1))&\xrightarrow{\mathrm{Gysin}}&H^{p+2}(X,L)\\
\cup\,\xi_Y^{\,n-1-p}\downarrow&&\downarrow\cup\,\xi_X^{\,n-1-p}\\
H^{2n-2-p}(Y,(L|Y)(n-2-p))&\xrightarrow{\mathrm{Gysin}}&
H^{2n-p}(X,L(n-1-p)).
\end{array}
\]
此图中, 当 $2n-2-p\ge n$, 即 $p\le n-2$ 时, 下行为同构 (7.2 (ii)).

关于这些猜想及其推论的讨论, 请读者参见 Kleiman 的报告 \emph{Algebraic cycles and the Weil conjectures}, 收录于 \emph{Dix exposés sur la cohomologie des schémas}.

\section*{8. 爆破簇}

本节给出 SGA 6 VII 3.7 和 4.8 中关于爆破概形的 $K$ 群及“代数闭链类群”的计算的上同调变体. 假设和记号采用第 4 节开头的那些 (4.1 至 4.3). 在这些条件下, 对“基本类”态射
\[
\gamma:j_*(\mu_{Y'}^{\otimes -1})[-2]\to A_{X'}
\]
施以函子 $Rf_*$ 后得到态射
\[
Rf_*(\gamma):i_*Rg_*(\mu_{Y'}^{\otimes -1})[-2]\to Rf_*(A_{X'}).
\]

\begin{statement}{定理 8.1}
$D^+(X,A)$ 中的态射
\[
A_X\oplus i_*Rg_*(\mu_{Y'}^{\otimes -1})[-2]\to
Rf_*(A_{X'})\oplus i_*(\mu_Y^{\otimes -d})[-2d]
\tag{8.1.1}
\]
由矩阵
\[
\begin{pmatrix}
\operatorname{adj}&Rf_*(\gamma)\\
0&i_*(\operatorname{trace})
\end{pmatrix}
\]
是同构.
\end{statement}

\emph{证明.} 由纤维函子族的保守性, 只需在 $X$ 的每个点 $x$ 上验证断言. 在 $U=X-Y$ 上显然成立. 对于 $x\in Y$, 利用固有基变换, 可归约为证明下列引理, 其中 $r=d-1$.

\begin{statement}{引理 8.2}
设 $k$ 为可分闭域, $r\ge1$ 为整数, $P=\mathbb P_k^r$. 则:
\begin{enumerate}[label=\alph*)]
\item 典范映射 $A\to H^0(P,A)$ 为双射;
\item $H^1(P,A)=H^{2r-1}(P,A)=0$;
\item 对 $0\le p\le2r-2$, 乘以 $\xi=c_1(\mathcal O_P(1))\in H^2(P,\mu)$ 的映射
\[
H^p(P,\mu_P^{\otimes -1})\to H^{p+2}(P,A)
\]
为同构;
\item 迹态射
\[
H^{2r}(P,\mu_P^{\otimes -1})\to H^0(k,\mu_k^{\otimes -(r+1)})
\]
为双射.
\end{enumerate}
\end{statement}

本引理是 2.2.2 的直接推论.

为陈述以下推论, 约定对每个 $X$-概形 $h:T\to X$、$T$ 的每个闭子集 $Z$ 和每个 $A_X$-模复形 $M$, 记
\[
H^*(T,M)=H^*(T,h^*(M)),\qquad
H^*_Z(T,M)=H^*_Z(T,h^*(M)).
\]

\begin{statement}{推论 8.3}
对每个 $A_X$-模复形 $M$ 和每个 $p\in\mathbb Z$, 映射
\[
H^p(X,M)\oplus H^{p-2}(Y',M\otimes\mu^{\otimes -1})
\to
H^p(X',M)\oplus H^{p-2d}(Y,M\otimes\mu^{\otimes -d})
\tag{8.3.1}
\]
与
\[
H_Y^p(X,M)\oplus H^{p-2}(Y',M\otimes\mu^{\otimes -1})
\to
H_{Y'}^p(X',M)\oplus H^{p-2d}(Y,M\otimes\mu^{\otimes -d})
\tag{8.3.2}
\]
由矩阵
\[
\begin{pmatrix}
f^*&j_*\\
0&g_*
\end{pmatrix}
\]
为双射.
\end{statement}

\emph{证明.} 暂用 $u$ 表示态射 (8.1.1). 由投影公式, 箭头 (8.3.1) 和 (8.3.2) 分别是态射 $R\Gamma(X,u\otimes\operatorname{id}_M)$ 及 $R\Gamma_Y(X,u\otimes\operatorname{id}_M)$ 在第 $p$ 个上同调对象上诱导的映射. 断言随即成立.

\subsection*{8.4}

现在明确 (8.3.2) 的逆同构, 为此需要若干引理.

\begin{statement}{引理 8.4.1}
设 $S$ 为概形, $E$ 为秩恒为 $d$ 的局部自由 $\mathcal O_S$-模, $p:P=\mathbb P_S(E)\to S$ 为相关联的射影丛. 设 $F$ 为由下列典范正合列定义的 $\mathcal O_P$-模:
\[
0\to F\to p^*(E)\to\mathcal O_P(1)\to0.
\]
则
\[
p_*(c_{d-1}(\check F))=1.
\]
\end{statement}

\emph{证明.} 用 $c_i$ 表示 $E$ 的 Chern 类, 令 $\xi=c_1(\mathcal O_P(1))$. 从
\[
p^*(c(E))=c(F)(1+\xi)
\]
得到
\[
c(F)=p^*(c(E))(1+\xi)^{-1},
\]
所以
\[
c_{d-1}(F)=\sum_{0\le i\le d-1}(-1)^{d-1-i}c_i\,\xi^{d-1-i}.
\tag{8.4.2}
\]
而 $p_*$ 将次数降低 $2d-2$, 且 $p_*(\xi^{d-1})=1$ (2.2.2). 对 (8.4.2) 两边施以 $p_*$ 并使用投影公式, 即得引理.

\begin{statement}{引理 8.4.3}
用 $F$ 表示由典范正合列
\[
0\to F\to g^*(N)\to\mathcal O_{Y'}(1)\to0.
\]
定义的局部自由 $\mathcal O_{Y'}$-模. 设 $M$ 为在 $Y$ 的某个邻域内局部常值且可构造的 $A_X$-模, $p\ge0$ 为整数. 对每个 $y\in H^p(Y,M)$, 有
\[
f^*(i_*(y))=j_*(g^*(y)c_{d-1}(\check F))
\]
等式在 $H_{Y'}^{p+2d}(X',M\otimes\mu^{\otimes d})$ 中成立.
\end{statement}

\emph{证明.} 令
\[
z=f^*(i_*(y))-j_*(g^*(y)c_{d-1}(\check F)).
\]
则
\[
j^*(z)=g^*i^*i_*(y)-j^*j_*(g^*(y)c_{d-1}(\check F)),
\]
由 (4.1) 及 (SGA 4 XVIII) 的线性引理,
\[
j^*(z)=g^*(y)\bigl(c_d(g^*\check N)-c_{d-1}(\check F)c_1(\mathcal O_{Y'}(-1))\bigr)=0.
\tag{8.4.4}
\]
记 $U=X-Y$, $U'=X'-Y'$, 有显然的正合交换图
\[
\begin{array}{ccccc}
H^{p+2d-1}(U')&\xrightarrow{\partial}&H^{p+2d}_{Y'}(X')&\xrightarrow{j^*}&H^{p+2d}(Y')\\
\uparrow\wr&&\uparrow f^*&&\uparrow g^*\\
H^{p+2d-1}(U)&\xrightarrow{\partial}&H_Y^{p+2d}(X)&\xrightarrow{i^*}&H^{p+2d}(Y),
\end{array}
\tag{8.4.5}
\]
其中记号作了显然的简化. 结合 (8.4.5) 及 (8.4.4), 可见 $z$ 形如
\[
z=f^*(t),
\]
其中 $t\in H_Y^{p+2d}(X,M\otimes\mu^{\otimes d})$. 由 (4.2), $t=f_*(z)$, 所以只需证明 $f_*(z)=0$. 而
\[
f_*(z)=f_*f^*i_*(y)-f_*j_*(g^*(y)c_{d-1}(\check F)),
\]
再用 (4.2), 得到
\[
f_*(z)=i_*(y)-i_*g_*(g^*(y)c_{d-1}(\check F)).
\]
由 $g$ 的投影公式及 (8.4.2) 即得结论.

\begin{statement}{命题 8.5}
设 $M$ 为在 $Y$ 的某个邻域内局部常值且可构造的 $A_X$-模.
\begin{enumerate}[label=\alph*)]
\item 设
\[
\gamma:H_{Y'}^*(X',M\otimes\mu^{\otimes d})\to
H^*(Y',M\otimes\mu^{\otimes(d-1)})(-2)
\]
为满足下述关系的唯一映射:
\[
f^*\circ f_* = j_*\circ\gamma+\operatorname{id}.
\tag{8.5.1}
\]
则
\[
g_*\circ\gamma=0,
\tag{8.5.2}
\]
且 (8.3.2) 的逆矩阵是
\[
\begin{pmatrix}
f_*&-i_*\\
-\gamma&g^*(\cdot)c_{d-1}(\check F)
\end{pmatrix}.
\tag{8.5.3}
\]
\item 对每个整数 $p\ge0$, 序列
\[
0\to H^{p-2d}(Y,M(-d))
\xrightarrow{\lambda}
H^{p-2}(Y',M(-1))\oplus H^p(X,M)
\xrightarrow{\mu}
H^p(X',M)\to0,
\]
其中
\[
\lambda(x)=\bigl(g^*(x)c_{d-1}(\check F),-i_*(x)\bigr),
\qquad
\mu(x,y)=j_*(x)+f^*(y),
\]
正合. 此外, 映射 $\lambda$ 有左逆
\[
\lambda'(x,y)=g_*(x).
\]
\end{enumerate}
\end{statement}

\emph{证明.} 先证明 a). 由 4.2 ($f_*f^*=\operatorname{id}$), 为证明
\[
\begin{pmatrix}
f_*&-i_*\\
-\gamma&g^*(\cdot)c_{d-1}(\check F)
\end{pmatrix}
\circ
\begin{pmatrix}
f^*&j_*\\
0&g_*
\end{pmatrix}
=\operatorname{id},
\]
只需证明
\[
\gamma j_*+\operatorname{id}=g^*(g_*(\cdot)c_{d-1}(\check F)),
\tag{8.5.4}
\]
以及
\[
\gamma f^*=0.
\tag{8.5.5}
\]
由 (8.5.1) 得到
\[
f^*f_*f^*=j_*\gamma f^*+f^*,
\]
再由 (4.2), 得到 $j_*\gamma f^*=0$; 因为 $j_*$ 是同构, 所以得到 (8.5.5). 下面证明 (8.5.4). 再由 (8.5.1) 得到
\[
f^*f_*j_*=j_*\gamma j_*+j_*,
\]
因为 $f_*j_*=i_*g_*$, 再用 (8.4.3), 得到
\[
j_*g^*(g_*(\cdot)c_{d-1}(\check F))=j_*\gamma j_*+j_*.
\]
左乘 $j_*$ 的逆即得 (8.5.4). 最后, 按另一顺序将两个矩阵相乘, 利用其积为恒等矩阵, 即得 (8.5.2).

下面证明 b). 由 8.4.3 及 8.4.1 分别得到 $\mu\circ\lambda=0$ 和 $\lambda'\circ\lambda=\operatorname{id}$. 此外, 由 (8.3), $\mu$ 为满射. 只需证明 $\Ker(\mu)\subset\operatorname{Im}(\lambda)$. 而 $\mu\circ(\lambda\circ\lambda')=0$, $\lambda'\circ(\lambda\circ\lambda')=\operatorname{id}$, 所以 $\mu(x)=0$ 蕴含
\[
(\mu\oplus\lambda')(\lambda\circ\lambda'(x))=(\mu\oplus\lambda')(x),
\]
再由 8.3, 得到 $x=\lambda\circ\lambda'(x)$.

\subsection*{8.6}

现在尝试确定 $X'$ 的上同调的乘法结构. 对每个 $p\in\mathbb Z$ 及每个在 $Y$ 的某个邻域内局部常值且可构造的 $A_X$-模 $M$, 令
\[
\nabla^p(X',M)=H^{p-2}(Y',M(-1))\oplus H^p(X,M).
\]
对于 $p,q\in\mathbb Z$ 及两个这样的模 $M,N$, 定义合成律
\[
\cup:\nabla^p(X',M)\otimes_A\nabla^q(X',N)\to
\nabla^{p+q}(X',M\otimes_A N)
\tag{8.6.1}
\]
定义为
\[
(y_1',x_1)\cup(y_2',x_2)=
\bigl(-y_1'y_2'\xi+y_1'g^*i^*(x_2)+(-1)^{pq}y_2'g^*i^*(x_1),\,x_1x_2\bigr).
\tag{8.6.2}
\]
此合成律显然反交换; 一个简单但繁琐的计算表明它满足结合律. 特别地,
\[
\nabla^*(X',A)=\bigoplus_{p\in\mathbb Z}\nabla^p(X',A)
\]
也带有分次交换 $A$-代数结构, 这里分次指 $\mathbb Z$ 分次, 且
\[
\nabla^*(X',M)=\bigoplus_{p\in\mathbb Z}\nabla^p(X',M)
\]
带有分次环 $\nabla^*(X',A)$ 上的分次模结构.

\begin{statement}{命题 8.6.3}
\begin{enumerate}[label=\alph*)]
\item 对每对在 $Y$ 的某个邻域内局部常值且可构造的 $A_X$-模 $(M,N)$, 下图
\[
\begin{tikzcd}
\nabla^*(X',M)\otimes\nabla^*(X',N)\arrow[r,"\mu\otimes\mu"]\arrow[d,"{(8.6.1)}"']&
H^*(X',M)\otimes H^*(X',N)\arrow[d,"\cup"]\\
\nabla^*(X',M\otimes N)\arrow[r,"\mu"]&
H^*(X',M\otimes N)
\end{tikzcd}
\]
交换. 特别地, 满射
\[
\mu_A=\mu:\nabla^*(X',A)\to H^*(X',A)
\]
为环态射, 且
\[
\mu_M=\mu:\nabla^*(X',M)\to H^*(X',M)
\]
为底环可变的模态射.
\item 若让 $\nabla^*(X',A)$ 通过映射
\[
\nabla^p(X',A)\otimes H^{q-2d}(Y,M(-d))
\to H^{p+q-2d}(Y,M(-d)),
\qquad
(y',x)\otimes y\mapsto i^*(x)y,
\]
作用于 $H^*(Y,M(-d))(-2d)$, 则 8.5 b) 中的映射 $\lambda$ 成为 $\nabla^*(X',A)$-模态射.
\end{enumerate}
\end{statement}

\emph{证明.} 先证明 a). 用 8.6.2 的记号, 需证明等式
\[
j_*\bigl(-y_1'y_2'\xi+y_1'g^*i^*(x_2)+(-1)^{pq}y_2'g^*i^*(x_1)\bigr)+f^*(x_1x_2)
\]
\[
=\bigl(j_*(y_1')+f^*(x_1)\bigr)\bigl(j_*(y_2')+f^*(x_2)\bigr).
\]
这立即由下列各式得到:
\[
\begin{cases}
f^*(x_1)f^*(x_2)=f^*(x_1x_2),&\text{由 }f^*\text{ 的乘法性},\\
j_*(-y_1'y_2'\xi)=j_*(y_1')j_*(y_2'),&\text{(4.5)},\\
j_*(y'g^*i^*(x))=j_*(y'j^*f^*(x))=j_*(y')f^*(x),&\text{投影公式}.
\end{cases}
\]
下面证明 b). 待证等式
\[
\lambda(i^*(x)\cdot y)=(y',x)\cdot\lambda(y),
\]
分解为下列两个等式:
\[
g^*(i^*(x)\cdot y)c_{d-1}(\check F)
=
-y'g^*(y)c_{d-1}(\check F)\xi
+y'g^*i^*(-i_*(y))
+g^*(i^*x)g^*(y)c_{d-1}(\check F),
\]
与
\[
-i_*(i^*(x)y)=-x\,i_*(y).
\]
第二个就是 $i$ 的投影公式. 第一个立即由
\[
-\xi\,c_{d-1}(\check F)=g^*[c_d(\check N)]
\quad\text{(8.4.3 的正合列)}
\]
以及
\[
i^*i_*(y)=y\,c_d(\check N)\qquad\text{(4.1)}.
\]

\section*{9. 爆破簇的 Chow 环及 Chow 环中的自交公式}

本节给出 Mumford 对 Chow 环中与 4.1 相应的公式的证明. 然后说明如何由它得到“关键公式” (8.4.3), 并进而确定爆破簇的 Chow 环的结构.

沿用 3.9 的约定, 用 $C^*(X)$ 表示域 $k$ 上的拟射影光滑概形 $X$ 的 Chow 环. 所用 Chern 类理论为 [1] 的理论. 以下引理将多次使用.

\begin{statement}{引理 9.1}
设 $S$ 为可分闭域 $k$ 上的拟射影光滑概形, $E$ 为秩恒为 $r+1$ 的局部自由 $\mathcal O_S$-模, $p:P=\mathbb P_S(E)\to S$ 为相关联的射影丛. 设 $F$ 为典范正合列
\[
0\to F\to p^*(E)\to\mathcal O_P(1)\to0,
\tag{9.1.0}
\]
定义的 $\mathcal O_P$-模, 并令 $\xi=c_1(\mathcal O_P(1))$. 则:
\begin{enumerate}[label=(\roman*)]
\item 对每个 $y\in C^*(S)$,
\[
y=p_*\bigl(p^*(y)c_r(\check F)\bigr);
\tag{9.1.1}
\]
\item 假设 $E=N\oplus\mathcal O_S$, 其中 $N$ 为秩 $r$ 的局部自由 $\mathcal O_S$-模, 于是 $P$ 等同于 $V(N)$ 的射影完备化 $\widehat N$. 设
\[
s:S\to\mathbb P_S(E)
\]
为 $V(N)$ 的零截面与典范嵌入 $V(N)\hookrightarrow\widehat N$ 复合而得的截面. 则:
\begin{enumerate}[label=(\alph*)]
\item 对每个 $y\in C^*(S)$, 有
\[
s_*(y)=p^*(y)c_r(\check F)=p^*(y)\left(\sum_{i=0}^r p^*(c_i(\check N))\xi^{r-i}\right)
\quad\text{在 }C^*(\widehat N)\text{ 中},
\tag{9.1.2}
\]
\[
s^*s_*(y)=y\,c_r(\check N)\quad\text{在 }C^*(S)\text{ 中};
\tag{9.1.3}
\]
\item 对每个 $z\in C^*(\widehat N)$,
\[
s^*(z)=p_*(z\,c_r(\check F))\quad\text{在 }C^*(S)\text{ 中}.
\tag{9.1.4}
\]
\end{enumerate}
\end{enumerate}
\end{statement}

\emph{证明.} 先证明 (9.1.1). 根据 $p$ 的投影公式, 只需证明 $p_*(c_r(\check F))=1$. 为此只要照搬 8.4.1 的证明, 使用关系 ([CH] 4 (17)). (9.1.2) 的第一项与第三项相等, 就是 [1] 的引理 3. 沿用 (8.4.1) 的证明可见
\[
c_r(\check F)=\sum_{i=0}^r p^*(c_i(\check N))\xi^{r-i},
\tag{9.1.2 bis}
\]
这证明了 (9.1.2). 对 (9.1.2) 首尾两项施以 $s^*$, 并利用 $p\circ s=\operatorname{id}$, 可见 (9.1.3) 将由
\[
s^*\left(\sum_{i=0}^{r-1}p^*(c_i(\check N))\xi^{r-i}\right)=0,
\]
更可由
\[
s^*(\xi^i)=0\qquad(i>0).
\tag{9.1.5}
\]
得到. 下面证明 (9.1.5). 若 $\alpha:V(N)\hookrightarrow\widehat N$ 表示典范嵌入, 则显然 $\alpha^*(\mathcal O_P(1))\simeq\mathcal O_{V(N)}$, 因而由 Chern 类的函子性,
\[
\alpha^*(\xi)=\alpha^*(c_1(\mathcal O_P(1)))=c_1(\mathcal O_{V(N)})=0
\]
根据 $s$ 的定义立即得到断言. 下面证明 (9.1.4). 将使用等式
\[
\xi\,c_r(\check F)=c_{r+1}(p^*(\check N)\oplus\mathcal O'_P)=0,
\tag{9.1.6}
\]
它直接来自正合列 (9.1.0). 元素 $z$ 有唯一分解
\[
z=\sum_{i=0}^r p^*(a_i)\xi^i,\qquad a_i\in C^*(S),
\]
([CH] 4 定理 1). 由 (9.1.5) 得到 $s^*(z)=a_0$. 另一方面, (9.1.6) 表明 $z\,c_r(\check F)=p^*(a_0)c_r(\check F)$, 再由 (9.1.1) 即得结论.

现在陈述 Chow 环中的自交公式.

\begin{statement}{定理 9.2 (Mumford)}
设 $k$ 为可分闭域, $i:Y\hookrightarrow X$ 为拟射影光滑 $k$-概形之间纯余维 $d$ 的闭浸入. 若以 $N$ 表示 $i$ 的法层, 则对每个 $y\in C^*(Y)$, 有
\[
i^*i_*(y)=y\,c_d(\check N)
\]
等式在 $C^*(Y)$ 中成立.
\end{statement}

为证明此定理, 需要一些记号及辅助概形. 通过态射 $u:t\mapsto(t,0)$ 将 $X$ 嵌入 $Z=X\times_k\mathbb P^1_k$, 从而得到笛卡尔图
\[
\begin{tikzcd}
\widehat N \arrow[r,"j_1"] \arrow[d,"g_1"'] & Z' \arrow[d,"f_1"]\\
Y \arrow[r,"i_1"'] & Z ,
\end{tikzcd}
\]
其中 $Z'$ 为沿 $Y$ 爆破 $Z$ 所得的概形, $\widehat N=\mathbb P_Y(N\oplus\mathcal O_Y)$ 为向量丛 $V(N)$ 的射影完备化. 如同 9.1 的陈述, 设 $s:Y\hookrightarrow\widehat N$ 为 $V(N)$ 的零截面与典范嵌入 $V(N)\subset\widehat N$ 的复合.

嵌入 $Y\subset X$ 将 $Y\times\mathbb P^1$ 等同于 $X\times\mathbb P^1$ 的一个闭子概形. 设 $W'$ 为 $Y\times\mathbb P^1$ 关于 $f_1$ 的严格变换, 即沿 $Y=Y\times0$ 爆破 $Y\times\mathbb P^1$ 所得的 $Z'$ 的子概形. 显然 $W'\cap\widehat N$ 等同于 $\widehat N$ 的零截面; 又由于 $Y$ 在 $Y\times\mathbb P^1$ 中的余维为 $1$, 典范态射 $f_2:W'\to Y\times\mathbb P^1$ 是同构. 与 $W'$ 有关的其余记号由典范图 $(D_1)$ 固定.
\[
\begin{tikzcd}[column sep=3.2em,row sep=2.4em,cells={nodes={inner sep=1pt}}]
\text{(D}_1\text{)}\arrow[from=1-3,to=5-2,"g_1"']
& & \widehat N \arrow[rr,"j_1"] & & Z' \arrow[ddd,"f_1"] \\
& Y \arrow[ur,"s"] \arrow[rr,"\delta"'] \arrow[ddd,"\mathrm{id}"'] & &
W' \arrow[ur,"\beta"'] \arrow[ddd,"f_2"'] & \\
& & & & \\
& & & & X\times\mathbb P^1=Z \\
& Y \arrow[rr,"\gamma"'] & & Y\times\mathbb P^1 \arrow[ur,"\alpha"'] &
\end{tikzcd}
\]

类似地, 若以 $X'$ 表示 $Z'$ 的闭子概形, 它由沿 $Y$ 爆破 $X=X\times0$ 得到, 则显然 $\mathbb P(N)=X'\cap\widehat N$ 等同于 $\widehat N$ 的无穷远超平面 $H$. 与 $X'$ 有关的其余记号由典范图 $(D_2)$ 固定.
\[
\begin{tikzcd}[column sep=3.2em,row sep=2.4em,cells={nodes={inner sep=1pt}}]
\text{(D}_2\text{)}\arrow[from=1-3,to=5-2,"g_1"']
& & \widehat N \arrow[rr,"j_1"] & & Z' \arrow[ddd,"f_1"] \\
& H \arrow[ur,"k"] \arrow[rr,"j"'] \arrow[ddd,"g"'] & &
X' \arrow[ur,"v"'] \arrow[ddd,"f"'] & \\
& & & & \\
& & & & X\times\mathbb P^1=Z \\
& Y \arrow[rr,"i"'] & & X \arrow[ur,"u"'] &
\end{tikzcd}
\]

下面用引理给出随后会用到的一些关系. 用 $E$ 表示由典范正合列
\[
0\to E\to g_1^*(N\oplus\mathcal O_Y)\to\mathcal O_{\widehat N}(1)\to0,
\tag{9.3.0}
\]
定义的局部自由 $\mathcal O_{\widehat N}$-模, 并令
\[
\bar\xi=c_1(\mathcal O_{\widehat N}(1)).
\tag{9.3.1}
\]

\begin{statement}{引理 9.3}
沿用上述记号, 有
\[
j_1^*(\beta_*(1))=c_d(\check E)\quad\text{在 }C^*(\widehat N)\text{ 中},
\tag{9.3.2}
\]
\[
j_1^*(v_*(1))=\bar\xi\quad\text{在 }C^*(\widehat N)\text{ 中},
\tag{9.3.3}
\]
\[
f_1^*u_*(1)=(j_1)_*(1)+v_*(1)\quad\text{在 }C^*(Z')\text{ 中}.
\tag{9.3.4}
\]
\end{statement}

\emph{证明.} 所考虑的全部 $k$-概形都光滑, 故由维数立即可见, 态射 $j_1$ 与 $\beta$ (分别地 $j_1$ 与 $v$) 横截. 由 ([CH] 4 命题 2 的推论 1), 得到
\[
\begin{cases}
j_1^*\beta_*(1)=s_*(1),\\
j_1^*v_*(1)=k_*(1).
\end{cases}
\]
等式 (9.3.2) 和 (9.3.3) 分别来自 (9.1.2) 及 Chern 类的归一化 ([CH] 4--18). 下面证明 (9.3.4). 由于 $Z'$ 的代数闭链 $X'$ 是 $X\times0$ 关于 $f_1$ 的严格变换, 预先有形如
\[
f_1^*u_*(1)=(j_1)_*g_1^*(\mu)+v_*(1),
\tag{9.3.5}
\]
的等式, 其中 $\mu\in\Gamma(Y,\mathbb N)$; 只需证明 $\mu=1$. 对 (9.3.5) 两边施以 $j_1^*$, 得到 (参见 $(D_2)$)
\[
g_1^*i^*u^*u_*(1)=j_1^*(j_1)_*g_1^*(\mu)+j_1^*v_*(1).
\tag{9.3.6}
\]
由 (9.1.3) 的一个极特殊情形, 或直接验证, 可得 $u^*u_*(1)=0$. 利用 (9.3.3), 从 (9.3.6) 得到
\[
0=j_1^*(j_1)_*g^*(\mu)+\bar\xi.
\tag{9.3.7}
\]
限制到 $Y$ 的一个连通分支的邻域后, 可假设 $Y$ 连通, 从而 (9.3.7) 成为
\[
0=\mu\,j_1^*(j_1)_*(1)+\bar\xi,\qquad \mu\in\mathbb N.
\tag{9.3.7 bis}
\]
由于 $j_1^*(j_1)_*(1)$ 为 $1$ 次齐次元素, 它形如
\[
j_1^*(j_1)_*(1)=g_1^*(\rho)+\nu\bar\xi,
\]
其中 $\rho\in C^1(Y)$, $\nu\in\mathbb Z$. 将此表达式代入 (9.3.7 bis), 立即得到 $\rho=0$, $\mu\nu=-1$, 其中 $\mu\in\mathbb N$; 所以 $\mu=1$, $\nu=-1$, 完成 (9.3.4) 的证明.

9.2 的证明的关键在于下述引理的证明. 如果像在上同调中那样, 已知对每个 $y\in C^*(\widehat N)$ 都有
\[
j_1^*(j_1)_*(y)=-y\bar\xi
\]
那么引理就是显然的.

\begin{statement}{引理 9.4}
对每个 $y\in C^*(\widehat N)$, 有
\[
c_d(\check E)\,j_1^*(j_1)_*(y)=0.
\]
\end{statement}

\emph{证明.} 元素 $y$ 在 $C^*(\widehat N)$ 中具有唯一的下述分解:
\[
y=\sum_{0\le p\le d}g_1^*(a_p)\bar\xi^p,
\]
或由 (9.3.3),
\[
y=\sum_{0\le p\le d}g_1^*(a_p)\bigl(j_1^*v_*(1)\bigr)^p.
\tag{9.4.1}
\]
利用 $j_1$ 的投影公式 ([CH] 4 I 9) 及 (9.1.6), 得到
\[
c_d(\check E)\,j_1^*(j_1)_*(y)=c_d(E)\,j_1^*(j_1)_*g_1^*(a_0);
\]
换言之, 归约为对形如 $y=g_1^*(a)$, $a\in C^*(Y)$ 的元素证明 9.4. 此时 $j_1^*(j_1)_*g_1^*(a)$ 在 $C^*(\widehat N)$ 中又具有唯一的下述分解:
\[
j_1^*(j_1)_*g_1^*(a)=\sum_{0\le p\le d}g_1^*(b_p)\bar\xi^p,
\tag{9.4.2}
\]
再次利用 (9.1.6), 只需证明 $g_1^*(b_0)c_d(\check E)=0$. 实际上, 将证明更强的断言 $b_0=0$. 由 (9.3.3), 可将 (9.4.2) 改写为
\[
j_1^*(j_1)_*g_1^*(a)=g_1^*(b_0)+z\,j_1^*v_*(1),
\tag{9.4.3}
\]
其中 $z\in C^*(\widehat N)$. 对 (9.4.3) 两边施以 $(j_1)_*$, 利用 $j_1$ 的投影公式, 得到
\[
(j_1)_*g_1^*(b_0)=\bigl[(j_1)_*g_1^*(a)\bigr](j_1)_*(1)-v_*(1)(j_1)_*(z).
\tag{9.4.4}
\]
将从 (9.3.4) 得到的 $(j_1)_*(1)$ 的表达式代入 (9.4.4), 得到
\[
(j_1)_*g_1^*(b_0)=
(j_1)_*g_1^*(a)f_1^*u_*(1)-v_*(1)\bigl[(j_1)_*g_1^*(a)+(j_1)_*(z)\bigr].
\tag{9.4.5}
\]
由于 $X\times0$ 与 $X\times1$ 在 $X\times\mathbb P^1$ 中线性等价, 有
\[
(j_1)_*g_1^*(a)f_1^*u_*(1)=0.
\tag{9.4.6}
\]
与 (9.4.5) 比较, 得到
\[
(j_1)_*g_1^*(b_0)=-v_*(1)\bigl[(j_1)_*g_1^*(a)+(j_1)_*(z)\bigr].
\tag{9.4.7}
\]
此外, 显然 $W'\cap X'=\emptyset$: 两者至多在 $Y$ 上方相交, 而在该处它们分别给出 $\widehat N$ 的零截面和无穷远超平面. 所以 $\beta^*v_*(1)=0$, 由 (9.4.7) 得到
\[
\beta^*(j_1)_*g_1^*(b_0)=0.
\tag{9.4.8}
\]
而如 9.3.2 的证明中所述, $\beta$ 与 $j_1$ 横截; 故由 (9.4.8) 得到
\[
0=\delta_*s^*g_1^*(b_0)=\delta_*(b_0).
\tag{9.4.9}
\]
最后, 对 (9.4.9) 施以 $(f_2)_*$, 得到 $\gamma_*(b_0)=0$, 从而 $b_0=0$, 因为 $\gamma_*$ 为单射 ($\gamma$ 有显然的收缩映射).

\subsection*{9.5}

最后说明上述结果如何给出定理. 设 $p:X\times\mathbb P^1_k\to X$ 为第一投影, 从而 $i=p\circ i_1$. 若 $y\in C^*(Y)$, 则显然
\[
i^*i_*(y)=(i_1)^*(\bar x),\qquad \bar x=p^*i_*(y).
\tag{9.5.1}
\]
由 (9.1.1), 得到
\[
i^*i_*(y)=(g_1)_*\bigl[g_1^*i_1^*(\bar x)c_d(\check E)\bigr],
\]
所以
\[
i^*i_*(y)=(g_1)_*\bigl[j_1^*f_1^*(\bar x)\cdot c_d(\check E)\bigr].
\tag{9.5.2}
\]
为证明定理, 由线性性可假设 $y$ 是与 $Y$ 的某个子簇 $T$ 对应的不可约代数闭链, 于是 $\bar x$ 对应于 $Z=X\times\mathbb P^1$ 的代数闭链 $T\times\mathbb P^1$. 用 $\theta$ 表示 $T\times\mathbb P^1$ 关于 $f_1$ 的严格变换, 即 $T\times\mathbb P^1$ 在 $Z\setminus Y$ 上的逆像的闭包. 则
\[
f_1^*(\bar x)=\theta+(j_1)_*(w),\qquad w\in C^*(\widehat N).
\tag{9.5.3}
\]
将此表达式代入 (9.5.2), 由 9.4 得到
\[
i^*i_*(y)=(g_1)_*\bigl[j_1^*(\theta)c_d(\check E)\bigr],
\tag{9.5.4}
\]
再由 (9.1.4),
\[
i^*i_*(y)=s^*j_1^*(\theta).
\tag{9.5.5}
\]
而 $\theta$ 的底约化概形是沿 $T=(T\times\mathbb P^1)\cap Y$ 爆破 $T\times\mathbb P^1$ 所得的概形; 它沿 $T$ 在 $\widehat N$ 的零截面下的像与该零截面横截相交. 因此
\[
j_1^*(\theta)=s_*(y),
\]
结合 (9.5.5) 和 (9.1.3), 完成定理的证明.

现在转向引言中提及的“关键公式”. 所用记号大部分沿用第 4 节, 这里作了适当调整并加以回顾. 注意它与本节引入的记号相容, 唯一例外是第 4 节记为 $Y'$ 的概形在这里记为 $H$ (取自“超平面”). 两套记号将同时使用.

\begin{statement}{命题 9.6}
设 $k$ 为可分闭域, $i:Y\hookrightarrow X$ 为拟射影光滑 $k$-概形之间纯余维 $d$ 的闭浸入. 设 $f:X'\to X$ 为沿 $Y$ 爆破 $X$ 所得的 $X$-概形, 并设
\[
\begin{tikzcd}
Y'\arrow[r,"j"]\arrow[d,"g"']&X'\arrow[d,"f"]\\
Y\arrow[r,"i"']&X
\end{tikzcd}
\]
为由 $f$ 与 $i$ 得到的笛卡尔图. 若以 $F$ 表示由典范正合列
\[
0\to F\to g^*(N)\to\mathcal O_{Y'}(1)\to0,
\tag{9.6.1}
\]
定义的局部自由 $\mathcal O_{Y'}$-模, 则
\[
f^*i_*(y)=j_*(g^*(y)c_{d-1}(\check F))
\tag{9.6.2}
\]
对每个 $y\in C^*(Y)$, 该等式在 $C^*(X')$ 中成立.
\end{statement}

证明 9.6 之前, 先将图 $(D_1)$ 与 $(D_2)$ 的若干尚未建立的性质单独列成引理.

\begin{statement}{引理 9.7}
有
\[
c_d(\check E)=k_*(c_{d-1}(\check F))+(g_1)^*(c_d(\check N)),
\tag{9.7 i}
\]
\[
\beta_*f_2^*\gamma_*(y)=f_1^*\alpha_*\gamma_*(y)\qquad\text{对每个 }y\in C^*(Y),
\tag{9.7 ii}
\]
\[
f_1^*u_*(t)=(j_1)_*((i\circ g_1)^*(t))+v_*f^*(t)\qquad\text{对每个 }t\in C^*(X).
\tag{9.7 iii}
\]
\end{statement}

\emph{证明.} 先证明 (i). 令 $\xi=c_1(\mathcal O_{Y'}(1))=k^*(\bar\xi)$. 沿用 (8.4.1) 的证明, 得到
\[
c_d(\check E)=\sum_{p=0}^d g_1^*(c_p(\check N))\bar\xi^{\,d-p},
\tag{9.7.1}
\]
\[
c_{d-1}(\check F)=\sum_{q=0}^{d-1}g^*(c_q(\check N))\xi^{d-1-q}.
\tag{9.7.2}
\]
由 (9.7.2), 得到
\[
k_*[c_{d-1}(\check F)]
=\sum_{q=0}^{d-1}k_*k^*[g_1^*(c_q(\check N))\bar\xi^{\,d-1-q}],
\]
再由 $k$ 的投影公式及 $k_*(1)=\bar\xi$ 这一事实,
\[
k_*[c_{d-1}(\check F)]=\sum_{q=0}^{d-1}g_1^*(c_q(\check N))\bar\xi^{\,d-q}.
\tag{9.7.3}
\]
将 (9.7.1) 与 (9.7.3) 比较, 立即得到断言 (i). 下面证明 (ii). 可假设 $y$ 为 $Y$ 的不可约代数闭链, 相应子簇也记为 $y$. 此时代数闭链 $\gamma_*(y)=y\times0$ 与 $y\times1$ 有理等价, 所以只需验证代数闭链等式
\[
\beta_*f_2^*(y\times1)=f_1^*\alpha_*(y\times1).
\tag{9.7.4}
\]
由于 $(y\times1)\cap(y\times0)=\emptyset$, $f_1$ 在 $y\times1$ 的某个开邻域上方为同构, 从而得到 (9.7.4). 下面证明 (iii). 若 $p:X\times\mathbb P^1\to X$ 为第一投影, 则 (9.1.2) 给出
\[
u_*(t)=p^*(t)u_*(1),
\tag{9.7.5}
\]
与 (9.3.4) 比较得到
\[
f_1^*u_*(t)=(p\circ f_1)^*(t)\bigl[(j_1)_*(1)+v_*(1)\bigr].
\tag{9.7.6}
\]
再由 $j_1$ 与 $v$ 的投影公式, 得到
\[
f_1^*u_*(t)=(j_1)_*\bigl[(p\circ f_1\circ j_1)^*(t)\bigr]+v_*\bigl[(p\circ f_1\circ v)^*(t)\bigr],
\tag{9.7.7}
\]
这就是 (iii), 因为 $p\circ f_1\circ j_1=i\circ g_1$, $p\circ f_1\circ v=f$.

\subsection*{9.8}

现在证明 9.6. 首先证明
\[
v_*f^*i_*(y)=v_*j_*(g^*(y)c_{d-1}(\check F)),
\tag{9.8.1}
\]
然后证明 $v_*$ 是单射. 由 (9.7)(iii), 得到
\[
v_*f^*i_*(y)=f_1^*u_*i_*(y)-(j_1)_*[g_1^*i^*i_*(y)],
\]
所以由 (9.2),
\[
v_*f^*i_*(y)=f_1^*(i_1)_*(y)-(j_1)_*[g_1^*(y\,c_d(\check N))].
\tag{9.8.2}
\]
另一方面,
\[
v_*j_*(g^*(y)c_{d-1}(\check F))
=(j_1)_*\bigl[k_*k^*g_1^*(y)c_{d-1}(\check F)\bigr],
\]
或由 $k$ 的投影公式,
\[
v_*j_*(g^*(y)c_{d-1}(\check F))
=(j_1)_*\bigl[g_1^*(y)k_*(c_{d-1}(\check F))\bigr].
\]
再用 (9.7)(i), 得到
\[
v_*j_*(g^*(y)c_{d-1}(\check F))
=(j_1)_*[g_1^*(y)c_d(\check E)]-(j_1)_*[g_1^*(y\,c_d(\check N))].
\tag{9.8.3}
\]
比较 (9.8.2) 和 (9.8.3), 可见 (9.8.1) 等价于
\[
f_1^*(i_1)_*(y)=(j_1)_*[g_1^*(y)c_d(\check E)],
\tag{9.8.4}
\]
即将 $X$ 换成 $X\times\mathbb P^1$ 后与 (9.6.2) 对应的等式. 根据 (9.1.2),
\[
(j_1)_*[g_1^*(y)c_d(\check E)]=(j_1)_*s_*(y)=\beta_*\delta_*(y).
\]
而立即可见 (参见 $(D_1)$) $\delta_*=f_2^*\gamma_*$, 因此:
\[
(j_1)_*[g_1^*(y)c_d(\check E)]=\beta_*f_2^*\gamma_*(y).
\tag{9.8.5}
\]
此外, 显然
\[
f_1^*(i_1)_*(y)=f_1^*\alpha_*\gamma_*(y),
\tag{9.8.6}
\]
比较 (9.8.5)、(9.8.6) 和 (9.7)(ii), 得到 (9.8.4).

\begin{statement}{引理 9.8.7}
映射 $v_*$ 是单射.
\end{statement}

\emph{证明.} 设 $Z''$ 为沿 $W'$ 爆破 $Z'$ 所得的概形, 等价地, 也就是沿 $\widehat N\cup W'=f_1^{-1}(Y\times\mathbb P^1)$ 的爆破. $X$ 在 $Z''$ 中的严格变换 $X''$ 也就是 $X'$ 在 $Z''$ 中的严格变换. 记号由下列典范图确定:
\[
\begin{tikzcd}
X''\arrow[r,"v''"]\arrow[d,"f'"',"\wr"]&Z''\arrow[d,"f'_1"]\\
X'\arrow[r,"v"']&Z'.
\end{tikzcd}
\]
由于 $X'\cap W'=\emptyset$, 态射 $f'_1$ 在 $X'$ 的某个开邻域上方是同构, 特别地 $f'$ 是同构. 此外, 因为 $f'_1$ 与 $v$ 横截, 有
\[
(f'_1)^*\circ v_*=v''_*\circ(f')^*,
\tag{9.8.8}
\]
所以只需证明 $v''_*$ 为单射. 而根据爆破的交换性 (例如参见 Hironaka: \emph{Smoothing of algebraic cycles}, Amer. Journal of Math. (1968), 引理 4.1), $Z''$ 还可如下得到: 先沿 $Y\times\mathbb P^1$ 爆破 $X\times\mathbb P^1$, 得到概形 $Z'_1$, 再沿 $Y$ 的逆像爆破 $Z'_1$. 显然 $Z'_1=X'\times\mathbb P^1$, 投影 $Z'_1\to X\times\mathbb P^1$ 为乘积 $f\times\operatorname{id}_{\mathbb P^1}$; 因此有显然的笛卡尔图
\[
\begin{tikzcd}[column sep=huge,row sep=large]
H\arrow[r,hook,"{t\mapsto(t,0)}"]\arrow[d,"g"']&H\times\mathbb P^1\arrow[r,hook,"{j\times\operatorname{id}}"]\arrow[d,"{g\times\operatorname{id}}"']&X'\times\mathbb P^1=Z'_1\arrow[d,"{f\times\operatorname{id}}"]\\
Y\arrow[r,hook,"\gamma"]&Y\times\mathbb P^1\arrow[r,hook,"{i\times\operatorname{id}}"]&X\times\mathbb P^1
\end{tikzcd}
\]
$X\times0$ 关于 $f\times\operatorname{id}$ 的严格变换等同于 $X'\times0$, 它包含 $H=H\times0=(f\times\operatorname{id})^{-1}(Y)$. 若再沿 $H$ 爆破 $Z'_1$, 因为 $H$ 在 $X'$ 中的纯余维为 $1$, $X'$ 的严格变换典范同构于 $X'$ (这是已知的); 因此得到典范交换图
\[
\begin{tikzcd}[column sep=large,row sep=large]
&X'\arrow[r,hook,"v''"]\arrow[d,"f''"']&Z''\arrow[d,"f''_1"]\\
H\arrow[r,hook]&X'\arrow[r,hook,"v'"]&X'\times\mathbb P^1\\
&t\arrow[r,mapsto]&{(t,0)}
\end{tikzcd}
\]
回顾我们的目标是证明 $v''_*$ 为单射. 与 $(D_2)$ 比较, 可见这等价于在附加条件 $Y$ 在 $X$ 中余维为 $1$ 时证明 $v_*$ 为单射, 下面就证明这一点. 此时恒等式 $(f_1)_*\circ v_*=u_*\circ f_*$ 给出结论, 因为 $f_*$ 为同构, $u_*$ 为可裂单态射 ($u$ 有收缩映射).

现在陈述对应于 (8.5 b) 的定理.

\begin{statement}{定理 9.9}
对每个整数 $p\ge0$, 序列
\[
0\to C^{p-d}(Y)\xrightarrow{\lambda} C^{p-1}(Y')\oplus C^p(X)\xrightarrow{\mu} C^p(X')\to0
\]
正合, 其中 $\lambda$ 和 $\mu$ 为映射
\[
\lambda:x\mapsto\bigl(g^*(x)c_{d-1}(\check F),-i_*(x)\bigr),
\qquad
\mu:(x,y)\mapsto j_*(x)+f^*(y)
\]
此外, 映射 $\lambda$ 有左逆
\[
\lambda':(x,y)\mapsto g_*(x).
\]
\end{statement}

\emph{证明.} 最后一个断言就是 (9.1.1). 同样, (9.6.2) 表明 $\mu\circ\lambda=0$. 先证明 $\mu$ 为满射. 将用到
\[
f_*f^*(z)=z\qquad(z\in C^*(X)),
\tag{9.9.1}
\]
它立即来自 $f$ 的投影公式及显然的事实 $f_*(1)=1$. 设 $t\in C^p(X')$. 由 (9.9.1),
\[
f_*(t-f^*f_*(t))=0.
\]
由于 $f$ 诱导从 $U'=X'-Y'$ 到 $U=X-Y$ 的同构, $t-f^*f_*(t)$ 在 $U'$ 上的限制为零, 故有等式 ([CH] 4 定理 2)
\[
t-f^*f_*(t)=j_*(a)
\qquad\text{其中 }a\in C^*(Y').
\]
现在证明中间项处的正合性. 设 $a\in C^{p-1}(Y')$, $b\in C^p(X)$ 满足
\[
j_*(a)+f^*(b)=0.
\tag{9.9.2}
\]
对两边施以 $f_*$, 得到
\[
f_*j_*(a)+f_*f^*(b)=0,
\]
再由 (9.9.1),
\[
b=-i_*g_*(a).
\tag{9.9.3}
\]
令 $x=g_*(a)$, 只需证明
\[
a=g^*(x)c_{d-1}(\check F).
\tag{9.9.4}
\]
令
\[
z=a-g^*(x)c_{d-1}(\check F)\qquad(x=g_*(a)),
\]
并证明 $z=0$. 先证明 $j_*(z)=0$ 及 $g_*(z)=0$, 再说明满足这两个关系的元素 $z$ 必为零. 根据定义,
\[
j_*(z)=j_*(a)-j_*(g^*(x)\cdot c_{d-1}(\check F)),
\]
所以由 (9.6.2),
\[
j_*(z)=j_*(a)-f^*i_*(x).
\]
利用 (9.9.3) 将 $i_*(x)$ 换成 $-b$, 由 (9.9.2) 得到 $j_*(z)=0$. 根据 $x$ 的定义,
\[
g_*(z)=x-g_*(g^*(x)c_{d-1}(\check F))=0
\]
这是由定理的最后一个断言得到的. 还需解释为何 $j_*(z)=0$ 与 $g_*(z)=0$ 蕴含 $z=0$. 由于 $Y'=\mathbb P_Y(N)$, $z$ 有唯一表达式
\[
z=g^*(a_0)+g^*(a_1)\xi+\cdots+g^*(a_{d-1})\xi^{d-1},
\]
其中 $a_i\in C^{p-1-i}(Y)$. 由于 $g_*(z)=0$, $g$ 的投影公式和 $g_*$ 将次数降低 $d-1$ 这一事实, 给出 $a_{d-1}=0$. 另一方面, 由 (9.2),
\[
z\,\xi=j^*j_*(z)=0,
\]
所以
\[
g^*(a_0)\xi+\cdots+g^*(a_{d-1})\xi^{d-1}=0.
\]
因此各 $a_i$ 均为零, 从而 $z=0$.

\subsection*{9.10}

用 (9.9) 的记号, 映射
\[
\omega=\lambda\circ\lambda'
\]
是满足 $\operatorname{Im}(\omega)=\operatorname{Im}(\lambda)$ 的投影算子. 所以 $\mu$ 诱导同构
\[
\bar\mu:\Ker(\omega)\xrightarrow{\sim}C^p(X').
\]
而
\[
\Ker(\omega)=\Ker(g_*)\oplus C^p(X)\hookrightarrow C^{p-1}(Y')\oplus C^p(X).
\tag{9.10.1}
\]
事实上, 元素 $(x,y)\in C^{p-1}(Y')\oplus C^p(X)$ 属于 $\Ker(\omega)$ 当且仅当
\[
g^*g_*(x)\cdot c_{d-1}(\check F)=-i_*g_*(x)=0,
\]
由 (9.1.1), 也就是当且仅当 $g_*(x)=0$. 另一方面, 射影丛的 Chow 环的结构, 特别是关系
\[
g_*(\xi^i)=0\quad\text{当 }i<d-1,\qquad g_*(\xi^{d-1})=1,
\]
表明
\[
\Ker(g_*)=g^*C^{p-1}(Y)\oplus g^*C^{p-2}(Y)\xi\oplus\cdots\oplus g^*C^{p-d+1}(Y)\xi^{d-2}.
\]
换言之, $\bar\mu$ 可解释为同构
\[
C^{p-1}(Y)\oplus C^{p-2}(Y)\oplus\cdots\oplus C^{p-d+1}(Y)\oplus C^p(X)
\xrightarrow{\sim}C^p(X'),
\tag{9.10.2}
\]
其逆很容易给出.

如 Murre 向我指出, 当 $d=2$ 时, Samuel 已在 \emph{Relations d'équivalence}, Proceedings of the International Congress, Edinburgh 1958 中证明了 (9.10.2) 的断言.

当然, 在平展上同调或 $\ell$ 进上同调的框架中, 有与 (9.10.2) 相应的结论, 这次由 (8.5) 导出.

\subsection*{参考文献}

\begin{enumerate}[label={[\arabic*]}]
\item Grothendieck, \emph{La théorie des classes de Chern}, Bull. Soc. Math. de France \textbf{86} (1958), pp. 137--154.
\item Hodge, \emph{The theory and applications of harmonic integrals}, Cambridge University Press, 1952.
\item Hirzebruch, \emph{Topological Methods in Algebraic Geometry}, Springer Verlag.
\item Lefschetz, \emph{Algebraic Geometry}, Princeton University Press, 1953.
\item Mme Raynaud, \emph{Modules projectifs universels}, Inventiones Mathematicae vol. 6. fasc. 1 (1968), pp. 1--26.
\item[{[TF]}] Godement, \emph{Théorie des faisceaux}.
\item[{[CH]}] \emph{Anneaux de Chow et applications}, Séminaire Chevalley 1958.
\item[{[6]}] \emph{Théorème de Lefschetz et critères de dégénérescence}, Pub. I.H.E.S. n$^\circ$ 35.
\end{enumerate}


\newpage
\phantomsection\addcontentsline{toc}{section}{第 VIII 报告 --- 完美复形}
\section*{第 VIII 报告}

\begin{center}
\textbf{阿贝尔范畴与三角范畴的类群.}\\
\textbf{完美复形.}\\[1ex]
作者 A. Grothendieck, I. Bucur 整理
\end{center}

本报告简要概述随后各报告不可缺少的概念和结果, 尤其涉及伪凝聚复形和完美复形. 关于后两者更完整且证明更详细的叙述, 将在继本讨论班之后举行的概形的 Riemann--Roch 定理讨论班的前几篇报告中给出.\footnote{SGA 6 I}

\subsection*{1. 阿贝尔范畴的情形}

设 $\cC$ 为阿贝尔范畴, $\cD$ 为 $\cC$ 的全子范畴. 从 $\cD$ 到阿贝尔群 $G$ 的函数 $f$ 称为可加的, 如果对 $\cC$ 中每个正合列
\[
0\longrightarrow E'\longrightarrow E\longrightarrow E''\longrightarrow 0,
\]
只要 $E',E'',E$ 都是 $\cD$ 的对象, 就有
\[
f(E)=f(E')+f(E'').
\]
设 $f$ 为从 $\cD$ 到阿贝尔群 $G$ 的可加函数. 下列性质容易由定义得到:
\[
\alpha)\quad f(0)=0,\qquad
\beta)\quad f(X)=f(Y)\quad\text{若 }X\text{ 与 }Y\text{ 同构},
\]
\[
\gamma)\quad f(X\oplus Y)=f(X)+f(Y).
\]

\[
\delta)\quad
\text{设 }E=(E^i,d^i)\text{ 为 }\cD\text{ 中的有限复形}
\]
且对象 $B^{i+1}=\operatorname{Im}(d^i)$, $Z^i=\operatorname{Ker}(d^i)$ 及 $H^i(E)=Z^i/B^i$ 都在 $\cD$ 中. 则
\[
\sum_i (-1)^i f(E^i)=\sum_i (-1)^i f(H^i(E)).
\]

\[
\varepsilon)\quad
\text{若 }F\in\Ob \cD\text{ 有有限滤过}
\]
\[
0=F_0\subset F_1\subset\cdots\subset F_n=F,
\]
且假设对 $0<i\le n$, $F_i$ 与 $F_i/F_{i-1}$ 都是 $\cD$ 的对象, 则
\[
f(F)=\sum_{i=1}^n f(F_i/F_{i-1}).
\]

\[
\varphi)\quad
\text{若 }\cD\text{ 对无限直和稳定, 则 }f=0.
\]

为证明 $\delta)$, 使用正合列
\[
0\longrightarrow Z^i\longrightarrow E^i\longrightarrow B^{i+1}\longrightarrow 0,
\]
\[
0\longrightarrow B^i\longrightarrow Z^i\longrightarrow H^i(E)\longrightarrow 0.
\]
$\varepsilon)$ 的证明对 $n$ 归纳, 使用正合列
\[
0\longrightarrow F_{n-1}\longrightarrow F\longrightarrow F/F_{n-1}\longrightarrow 0.
\]
为证明 $\varphi)$, 设 $E\in\Ob\cD$, 并令
\[
F=\bigoplus_{n\in\mathbb Z^+}E_n,\qquad E_n=E,\quad n\in\mathbb Z^+.
\]
则 $E\oplus F\simeq F$, 应用 $\gamma)$ 得到
\[
f(F)=f(F)+f(E),
\]
所以 $f(E)=0$.

对每个阿贝尔群 $G$, 考虑由从 $\cD$ 到 $G$ 的可加函数组成的阿贝尔群 $A_{\cD}(G)$, 显然得到协变函子
\[
A_{\cD}:\Ab\longrightarrow\Ab.
\]
函子 $A_{\cD}$ 可表. 代表它的一对 $(K_{\cD}(\cC),\varphi)$ 可如下构造. 设 $L$ 为由 $\cD$ 的对象生成的自由阿贝尔群, $R$ 为 $L$ 中由形如 $E-E'-E''$ 的元素生成的子群, 每个这样的元素对应于一个正合列
\[
0\longrightarrow E'\longrightarrow E\longrightarrow E''\longrightarrow 0
\]
(在 $\cC$ 中). 令
\[
K_{\cD}(\cC)=L/R.
\]
映射 $\varphi$ 通过复合两个典范映射
\[
\Ob\cD\longrightarrow L\longrightarrow L/R.
\]
得到. 不会混淆时, 可以简写记号:
\[
K_{\cD}(\cC)=K(\cD).
\]

\subsection*{2. 三角范畴的情形}

设 $\cC$ 为三角范畴, $G$ 为阿贝尔群. 从 $\cC$ 到 $G$ 的函数 $f$ 称为可加的, 如果对每个特异三角形
\[
\begin{tikzcd}[column sep=small,row sep=small]
& Z \arrow[dl,"w"']\\
X \arrow[rr,"u"'] && Y \arrow[ul,"v"']
\end{tikzcd}
\]
有
\[
f(Y)=f(X)+f(Z).
\]

给定可加函数 $f$, 下列关系直接来自定义:
\[
\text{a) } f(0)=0,\qquad
\text{b) } f(X[n])=(-1)^n f(X),
\]
\[
\text{c) } f(X)=f(Y)\text{ 若 }X\text{ 与 }Y\text{ 同构},
\qquad
\text{d) } f(X\oplus Y)=f(X)+f(Y).
\]
事实上, 以下三角形都是特异三角形:
\[
\begin{tikzcd}[column sep=small,row sep=small]
&0 \arrow[dl]\\ X \arrow[rr,"1_X"']&&X\arrow[ul]
\end{tikzcd}
\qquad
\begin{tikzcd}[column sep=small,row sep=small]
&X[1]\arrow[dl]\\ X\arrow[rr]&&0\arrow[ul]
\end{tikzcd}
\]
\[
\begin{tikzcd}[column sep=small,row sep=small]
&0\arrow[dl]\\ X\arrow[rr,"u"']&&Y\arrow[ul]
\end{tikzcd}
\qquad
\begin{tikzcd}[column sep=small,row sep=small]
&Y[1]\arrow[dl]\\ X\oplus Y\arrow[rr]&&X\arrow[ul]
\end{tikzcd}
\]
($u$ 为同构). 构造最后一个三角形时, 使用了两个特异三角形之和仍为特异三角形这一事实.

如同阿贝尔范畴的情形, 显然得到函子
\[
A:\Ab\longrightarrow\Ab
\]
它将每个阿贝尔群 $G$ 送到 $\cC$ 上取值于 $G$ 的可加函数组成的阿贝尔群 $A(G)$. 此函子可表, 容易构造代表它的一对 $(K(\cC),\operatorname{cl}_{\cC})$. 具体地, 先考虑由 $\cC$ 的对象生成的自由阿贝尔群 $\mathbb Z^{(\Ob\cC)}$, 再取由形如 $X-Y-Z$ 的元素生成的子群 $R$, 其中
\[
\begin{tikzcd}[column sep=small,row sep=small]
&Z\arrow[dl]\\ Y\arrow[rr]&&X\arrow[ul]
\end{tikzcd}
\]
为特异三角形. 商群 $\mathbb Z^{(\Ob\cC)}/R$ 即 $K(\cC)$; $\operatorname{cl}_{\cC}$ 由两个典范映射
\[
\Ob\cC\longrightarrow \mathbb Z^{(\Ob\cC)}\longrightarrow K(\cC).
\]

\subsection*{3. 函子性质}

设 $\cC,\cC'$ 为两个三角范畴, $T:\cC\to\cC'$ 为从 $\cC$ 到 $\cC'$ 的正合函子, 即将特异三角形送到特异三角形的可加“分次”函子. 如前, 用
\[
A,A':\Ab\longrightarrow\Ab
\]
表示与范畴 $\cC$ 及 $\cC'$ 相关联的函子. 与 $T$ 复合给出函子间的态射 $A'\to A$, 因而得到同态
\[
K(T):K(\cC)\longrightarrow K(\cC')
\]
它由下图的交换性唯一确定:
\[
\begin{tikzcd}[column sep=large,row sep=large]
\Ob\cC \arrow[r,"T"] \arrow[d,"\operatorname{cl}_{\cC}"']&
\Ob\cC' \arrow[d,"\operatorname{cl}_{\cC'}"]\\
K(\cC) \arrow[r,"K(T)"']& K(\cC').
\end{tikzcd}
\]
因此, 若函子 $T,T_1:\cC\to\cC'$ 同构, 则 $K(T)=K(T_1)$. 特别地, 若 $T$ 为范畴等价, 则 $K(T)$ 是同构.

若 $\cC''$ 为第三个三角范畴, 且
\[
F:\cC\times\cC'\longrightarrow \cC''
\]
为关于每个变量都正合的双函子, 则显然得到双线性映射
\[
K(F):K(\cC)\times K(\cC')\longrightarrow K(\cC'')
\]
使得下列关系成立:
\[
K(F)(\operatorname{cl}_{\cC}(X),\operatorname{cl}_{\cC'}(X'))
=
\operatorname{cl}_{\cC''}(F(X,X')).
\]

\begin{statement}{命题 3.1}
设 $\cA$ 为三角范畴, $\cB$ 为 $\cA$ 的厚三角子范畴, $\cA/\cB$ 为 $\cA$ 关于 $\cB$ 的商范畴, $Q:\cA\to\cA/\cB$ 为典范取商函子, $i:\cB\to\cA$ 为包含函子. 则下列序列正合:
\[
K(\cB)\xrightarrow{K(i)}K(\cA)\xrightarrow{K(Q)}K(\cA/\cB)\longrightarrow 0.
\]
\end{statement}

\emph{证明.} 显然有自然同态
\[
\operatorname{Coker}(K(i))\xrightarrow{\alpha}K(\cA/\cB).
\]
另一方面, 有显然的映射
\[
\Ob(\cA/\cB)\xrightarrow{\gamma}\operatorname{Coker}(K(i)),
\]
因为 $\Ob(\cA/\cB)=\Ob(\cA)$. 设 $X,Y\in\Ob(\cA/\cB)$ 且它们同构. 要证明 $\gamma(X)=\gamma(Y)$. 范畴 $\cA/\cB$ 由右分式或左分式演算得到, 所用乘法系 $S$ 由包含在特异三角形 $(X',Y',Z',f,g,h)$ 中的态射 $f$ 组成, 其中 $Z'$ 为 $\cB$ 的对象. ($S$ 为饱和乘法系.) 因而存在 $X_1$ 及态射
\[
X_1\xrightarrow{s}X,\qquad X_1\xrightarrow{t}Y
\]
使得 $s,t\in S$. 因此只需证明 $\gamma(X)=\gamma(X_1)$, 而这来自存在形如
\[
(X_1,X,Z,s,\ldots),
\]
的特异三角形, 其中 $Z\in\Ob\cB$.

映射 $\gamma$ 为可加函数. 事实上, 若 $(X,Y,Z,f,g,h)$ 为 $\cA/\cB$ 的特异三角形, 需要证明
\[
\gamma(X)+\gamma(Z)=\gamma(Y).
\]
由于三角形 $(X,Y,Z,f,g,h)$ 同构于 $\cA$ 的一个特异三角形, 此关系随即成立. 因此得到
\[
\beta:K(\cA/\cB)\longrightarrow \operatorname{Coker}(K(i)),
\]
并且立即可验证 $\beta$ 为 $\alpha$ 的逆.

\subsection*{4. 与阿贝尔范畴情形的比较}

设 $\cC$ 为阿贝尔范畴, $D(\cC)$ 为 $\cC$ 的导出范畴. 如常, 用 $D^b(\cC)$ 表示 $D(\cC)$ 中由有界复形生成的全子范畴. 给 $\cC$ 的每个对象 $X$ 配以除零次分量等于 $X$ 外其余分量皆为零的复形, 得到同态
\[
K(\cC)\xrightarrow{\mu}K(D^b(\cC)).
\]
此同态实际上为同构. 事实上, 利用命题 3.1 可以构造 $\mu$ 的逆, 这里注意到
\[
D^b(\cC)\simeq K^{b,b}(\cC)/K^{b,\phi}(\cC),
\]
其中 $K^{b,\phi}(\cC)$ 为 $K^{b,b}(\cC)$ 中由无上同调复形组成的厚子范畴. 设 $X^\bullet=(X_i)\in\Ob(K^{b,b}(\cC))$. 若令
\[
f(X^\bullet)=\sum_i(-1)^i\operatorname{cl}_{\cC}(X_i),
\]
便得到从 $K^{b,b}(\cC)$ 到 $K(\cC)$ 的可加函数. 事实上, 首先注意若 $X^\bullet$ 与 $Y^\bullet$ 同伦, 则 $f(X^\bullet)=f(Y^\bullet)$, 因为 $f(X^\bullet)=f(H(X^\bullet))$. 因而有同态
\[
K(D^b(\cC))\xrightarrow{\nu}K(\cC),
\]
立即可验证它为 $\mu$ 的逆.

当 $\cC$ 过大时, 不变量 $K(\cC)$ 往往没有意义, 因而将它相对于 $\cC$ 的子范畴 $\cC_o$ 相对化是有用的. 设 $\cC_o$ 为 $\cC$ 的全子范畴. 一方面, 可以定义群
\[
K_o(\cC)=K_{\cC_o}(\cC).
\]
另一方面, 假设满足下列条件:
\[
(A)\qquad \cC_o\text{ 对有限直和稳定.}
\]
此时很容易定义三角范畴 $D_o^b(\cC)$: 考虑 $\cC$ 中上同调有界、且各分量在子范畴 $\cC_o$ 中的复形, 再关于这些复形的拟同构取商. 如前, 得到同态
\[
(*)\qquad K_o(\cC)\longrightarrow K(D_o^b(\cC)).
\]
假设子范畴 $\cC_o$ 还满足下列条件 (B):
\[
(B)\qquad
\text{若 }X,Y\in\Ob\cC_o\text{ 且 }X\xrightarrow{u}Y
\text{ 为满态射 (当然是在 }\cC\text{ 中)},
\]
则 $\operatorname{Ker}u\in\Ob\cC_o$. 如上论证, 可证明同态 $(*)$ 为同构.

典范包含函子
\[
D_o^b(\cC)\longrightarrow D^b(\cC)
\]
一般并非全忠实的. 但是:

\begin{statement}{命题 4.1}
假设子范畴 $\cC_o$ 满足条件 (A)、(B) 及下列条件 (C):
\[
(C)\qquad
\text{每个图 }
\begin{tikzcd}[column sep=large,row sep=large]
A \arrow[r,"u"]&B\\
&L \arrow[u,"v"']
\end{tikzcd}
\]
其中 $u$ 为满态射, $L\in\Ob\cC_o$, 都可补成交换图
\[
\begin{tikzcd}[column sep=large,row sep=large]
A \arrow[r,"u"]&B\\
L_1 \arrow[u] \arrow[r,"w"']&L \arrow[u,"v"']
\end{tikzcd}
\]
使得 $L_1\in\Ob\cC_o$, 且 $w$ 为满态射. 在这些条件下, 典范函子
\[
D_o^b(\cC)\longrightarrow D^b(\cC),\qquad
D_o^-(\cC)\longrightarrow D^-(\cC)
\]
全忠实. 此外, $D_o^b(\cC)$ (分别地 $D_o^-(\cC)$) 的本质像对直因子稳定.
\end{statement}

本命题的证明主要使用 Verdier 的一个判别准则 ([2], 4--2 定理), 将在 Riemann--Roch 讨论班的第一篇报告中给出 (SGA 6 I 2.7).

\begin{statement}{推论}
再假设子范畴 $\cC_o$ 对 $\Ker$ 和 $\Coker$ 稳定, 因而为阿贝尔范畴; 用 $D^b(\cC)_o$ 表示 $D^b(\cC)$ 的全三角子范畴, 它由上同调对象同构于 $\cC_o$ 的对象的那些 $D^b(\cC)$ 复形生成. 则包含函子
\[
D^b(\cC_o)\longrightarrow D^b(\cC)_o
\]
是两个范畴间的等价, 因而得到一连串同构
\[
K(\cC_o)\simeq K_o(\cC)\simeq K(D^b(\cC_o))\simeq K(D^b(\cC)_o).
\tag{*}
\]
\end{statement}

\emph{评注.} 假设子范畴 $\cC_o$ 满足条件 (A)、(B), 且 $\cC$ 的每个对象均有由 $\cC_o$ 的对象组成的有限左消解. 此时 $\cC_o$ 也满足条件 (C), 典范函子
\[
D_o^b(\cC)\longrightarrow D^b(\cC)
\]
实际上为范畴等价. 因此有同构
\[
K_{\cC_o}(\cC)\simeq K(\cC).
\]
\footnote{参见 (SGA 6 IV 1.6).}

\subsection*{5. 伪凝聚复形}

设 $\cC$ 为阿贝尔范畴, $\cC_o$ 为 $\cC$ 的全子范畴, 满足条件 (A)、(B)、(C). 用 $D_o^-(\cC)_{\coh}$ 表示 $D_o^-(\cC)$ 在典范函子
\[
D_o^-(\cC)\longrightarrow D(\cC).
\]
下的本质像. $D_o^-(\cC)_{\coh}$ 的对象称为 $\cC_o$-伪凝聚复形. 因而 $\cC$ 中的复形为 $\cC_o$-伪凝聚, 当且仅当它在 $D(\cC)$ 中同构于一个各分量都是 $\cC_o$ 的对象的上有界复形.

范畴 $D_o^-(\cC)_{\coh}$ 是 $D(\cC)$ 的三角子范畴, 意即 $D(\cC)$ 的每个有两个顶点在子范畴 $D_o^-(\cC)_{\coh}$ 中的特异三角形, 都同构于一个三个对象都在子范畴 $D_o^-(\cC)_{\coh}$ 中的三角形.

首先, 显然子范畴 $D_o^-(\cC)_{\coh}$ 在平移自同构 $T:D(\cC)\to D(\cC)$ 下稳定. 所以要证明 $D_o^-(\cC)_{\coh}$ 为 $D(\cC)$ 的三角子范畴, 只需验证: 若 $U^\bullet,V^\bullet$ 为 $D_o^-(\cC)_{\coh}$ 的对象, 且三角形
\[
\begin{tikzcd}[column sep=small,row sep=small]
& W^\bullet \arrow[dl,"w"']\\
U^\bullet \arrow[rr,"u"']&& V^\bullet \arrow[ul,"v"']
\end{tikzcd}
\]
在 $D(\cC)$ 的三角范畴结构下是特异的, 则 $W^\bullet$ 也是 $D_o^-(\cC)_{\coh}$ 的对象. 由假设, 存在复形 $X^\bullet,Y^\bullet$, 它们在 $D(\cC)$ 中分别同构于 $U^\bullet,V^\bullet$, 且各分量为 $\cC_o$ 的对象. 得到 $D(\cC)$ 中的交换方块:
\[
\begin{tikzcd}[column sep=large,row sep=large]
U^\bullet \arrow[r,"u"] \arrow[d,"\sim"']&V^\bullet \arrow[d,"\sim"]\\
X^\bullet \arrow[r,"f"']&Y^\bullet .
\end{tikzcd}
\]
不过, 将 $X^\bullet$ 换成另一同构复形后, 利用包含函子
\[
D_o^-(\cC)\longrightarrow D^-(\cC)
\]
全忠实这一事实 (命题 4.1), 可假设 $f$ 为真正的复形态射. 设 $Z^\bullet$ 为 $f$ 的映射柱, $g:Y^\bullet\to Z^\bullet$ 为典范态射. 由构造本身, $Z^\bullet$ 的各分量都是 $\cC_o$ 的对象 (条件 (A)). 另一方面,
\[
\begin{tikzcd}[column sep=small,row sep=small]
&Z^\bullet \arrow[dl,"h"']\\
X^\bullet \arrow[rr,"f"']&&Y^\bullet \arrow[ul,"g"']
\end{tikzcd}
\]
是 $D(\cC)$ 中的特异三角形, 因而 $Z^\bullet$ 同构于 $W^\bullet$.

所以范畴 $D_o^-(\cC)_{\coh}$ 带有三角范畴结构: $D_o^-(\cC)_{\coh}$ 中的三角形为特异三角形, 当且仅当它在三角范畴 $D(\cC)$ 中是特异三角形.

将 $D_o^-(\cC)_{\coh}$ 与 $D^b(\cC)_{\coh}$ 的交记为 $D_o^b(\cC)_{\coh}$.

$\cC$ 的对象 $M$ 称为伪凝聚的, 如果视它为仅集中于零次的复形时是伪凝聚的; 可证明, 其充要条件是它具有各分量为 $\cC_o$ 的对象的左消解.

\begin{statement}{命题 5.1}
复形 $X$ 为 $\cC_o$-伪凝聚, 当且仅当它同构于一个右有界、各分量均 $\cC_o$-伪凝聚的复形.
\end{statement}

本命题的证明将在 Riemann--Roch 讨论班的第一篇报告中给出 (SGA 6 I 2.7).

\subsection*{6. 完美复形}

范畴 $D(\cC)$ 的对象称为 $\cC_o$-完美复形, 如果它同构于 $D_o^b(\cC)$ 的对象, 即同构于一个各分量都在 $\cC_o$ 中的有界复形. 用 $D_o(\cC)_{\parf}$ 表示由完美复形生成的 $D(\cC)$ 的全子范畴.

所以范畴 $D_o(\cC)_{\parf}$ 正是 $D_o^b(\cC)$ 在包含函子
\[
D_o^b(\cC)\longrightarrow D(\cC).
\]
下的本质像. 容易看出 $D_o(\cC)_{\parf}$ 是 $D(\cC)$ 的三角子范畴, 因而带有三角范畴结构.

\subsection*{7. 一个重要特例}

设 $A$ 为环, $\cC$ 为 $A$-模范畴, 例如取左模. 用 $\cC_o$ 表示 $\cC$ 中由有限型投射 $A$-模生成的全子范畴. 子范畴 $\cC_o$ 满足条件 (A)、(B)、(C). 在这种情形下使用以下记号:
\[
D^-(A)_{\coh}=D_o^-(\cC)_{\coh},\qquad
D^b(A)_{\coh}=D_o^b(\cC)_{\coh},
\]
\[
D(A)_{\parf}=D_o(\cC)_{\parf}.
\]
$D(\cC)$ 中为 $\cC_o$-伪凝聚 (分别地 $\cC_o$-完美) 的复形, 简称伪凝聚 (分别地完美) 复形.

考虑以下各群是有用的:
\[
K_\bullet(A)=K(D^b(A)_{\coh}),
\]
\[
K^\bullet(A)=K(D(A)_{\parf}).
\]
至于群 $K^\bullet(A)$, 由命题 4.1 及其推论, 有
\[
K^\bullet(A)=K(D(A)_{\parf})\simeq K(D_o^b(\cC))\simeq K_{\cC_o}(\cC),
\]
可见它也就是 $A$ 上有限型投射模的经典类群.

在 $A$ 为诺特环的特殊情形下, 群 $K_\bullet(A)$ 也有经典描述. 为此, 设 $\cC_1$ 为有限型模的阿贝尔范畴. 则
\[
K_\bullet(A)\simeq K(\cC_1).
\]
事实上, 利用命题 4.1, 有
\[
K(\cC_1)\simeq K(D^b(\cC_1))\simeq K(D^b_{\cC_1}(\cC)),
\]
因为立即有 $D^b(\cC_1)\simeq D^b_{\cC_1}(\cC)$.

实际上, 若 $A$ 为诺特环, $L^\bullet\in D^-(A)$, 则由命题 5.1, $L^\bullet$ 伪凝聚, 当且仅当各 $H^i(L^\bullet)$ 为有限型.

对任意环 $A$, 有同态
\[
K^\bullet(A)\xrightarrow{\widetilde\gamma_A}K_\bullet(A),
\]
它来自包含函子
\[
(*)\qquad D(A)_{\parf}\longrightarrow D^b(A)_{\coh}.
\]
一般而言, 同态 $\widetilde\gamma_A$ 既非单射也非满射. 在 $A$ 为正则环的特殊情形下, 每个有限型 $A$-模 $M$ 均有有限投射维数, 再由 5.1, 包含函子 $(*)$ 实际上为两个范畴间的等价. 因而此时 $\widetilde\gamma_A$ 为同构.

\subsection*{8. 复形的 Tor}

设 $A$ 为不必交换的环, ${}_A\cC$ (分别地 $\cC_A$) 为左 (分别地右) $A$-模的阿贝尔范畴. 用 $K^-({}_A\cC)$ 表示 ${}_A\cC$ 中的上有界复形的三角范畴, 用 ${}_A\cD$ 表示 $K^-({}_A\cC)$ 中由各分量为平坦左 $A$-模的复形组成的全三角子范畴.

设 $X^\bullet$ 为 $K^-(\cC_A)$ 的对象, 考虑函子
\[
\Phi_{X^\bullet}=X^\bullet\otimes_A\,\cdot\;:\;K^-({}_A\cC)\longrightarrow K^-(\Ab).
\]
在这些条件下, 下列命题成立.

1) ${}_A\cD$ 的每个无上同调对象经 $\Phi_{X^\bullet}$ 映成 $K^-(\Ab)$ 的无上同调对象. 事实上, 若 $N^\bullet$ 为 ${}_A\cD$ 的对象, 且 $Z^i(N^\bullet)$、$B^i(N^\bullet)$ 与 $H^i(N^\bullet)$ 平坦, 则有经典同构
\[
\bigoplus_{p+q=n}H_p(X^\bullet)\otimes H_q(N^\bullet)
\longrightarrow H_n(X^\bullet\otimes_A N^\bullet).
\]

2) $K^-({}_A\cC)$ 的每个对象 $Y^\bullet$ 都是一个源为 ${}_A\cD$ 的对象的拟同构的目标. 为此只需取 Cartan--Eilenberg 消解. 因此可应用 [2] §2 第 2 节的定理 2, 得到 $\Phi_{X^\bullet}$ 有左导出函子:
\[
\mathbf L\Phi_{X^\bullet}:D^-({}_A\cC)\longrightarrow D^-(\Ab).
\]
对象 $\mathbf L\Phi_{X^\bullet}(Y^\bullet)$ 可在相差同构的意义下如下得到. 设 $P^\bullet$ 为 ${}_A\cD$ 的对象, $P^\bullet\to Y^\bullet$ 为拟同构. 则 $\mathbf L\Phi_{X^\bullet}(Y^\bullet)$ 同构于
\[
Q(X^\bullet\otimes_A P^\bullet),
\]
其中如常, $Q$ 表示典范函子
\[
K^-(\Ab)\longrightarrow D^-(\Ab).
\]
因此, 对 $D^-({}_A\cC)$ 的每个对象 $Y^\bullet$, 函子
\[
K^-(\cC_A)\longrightarrow D^-(\Ab),\qquad
X^\bullet\longmapsto \mathbf L\Phi_{X^\bullet}(Y^\bullet)
\]
正合且在无上同调复形上为零, 因而给出正合双函子
\[
D^-(\cC_A)\times D^-({}_A\cC)\longrightarrow D^-(\Ab),
\]
\[
(X^\bullet,Y^\bullet)\longmapsto X^\bullet\overset{\mathbf L}{\otimes}_A Y^\bullet.
\]
取同调, 定义 $\operatorname{Tor}_i^A(X^\bullet,Y^\bullet)$ (也参见 [1], 6.3). 在 $A$ 交换的特殊情形下, 实际上得到正合双函子
\[
D^-(\cC)\times D^-(\cC)\longrightarrow D^-(\cC),
\]
即导出范畴意义下的张量积 ($\cC={}_A\cC=\cC_A$).

\begin{statement}{定义}
若 $D(\cC_A)$ 的复形 $X$ 属于 $D^-$, 且存在 $n_o\in\mathbb Z^+$ 使得对每个 $Y\in\Ob({}_A\cC)$ 及每个 $n>n_o$, 都有
\[
\operatorname{Tor}_n^A(X,Y)=0
\]
则称 $X$ 有有限 Tor 维数.
\end{statement}

将利用 Tor 维数的概念来刻画完美复形.

\begin{statement}{命题 8.1}
设 $A$ 为不必交换的环, $L^\bullet$ 为 $A$-模复形. 下列条件等价:
\[
\text{a) } L^\bullet\text{ 为完美复形;}
\]
\[
\text{b) } L^\bullet\text{ 伪凝聚、上同调有界且有有限 Tor 维数.}\footnote{参见 SGA 6 I 5.8.1.}
\]
\end{statement}

证明使用以下引理.

\begin{statement}{引理}
(Bourbaki, Alg. Comm. I, §2, 习题 15). 每个平坦且有限表现的 $A$-模 $E$ 都是投射模.
\end{statement}

\emph{引理的证明.} 设 $u:F\to F'$ 为满态射. 需要证明
\[
\Hom_A(E,F)\longrightarrow\Hom_A(E,F')
\]
为满态射. 只需证明对每个可除群 $G$,
\[
\Hom_{\mathbb Z}(\Hom_A(E,F'),G)\longrightarrow
\Hom_{\mathbb Z}(\Hom_A(E,F),G)
\]
为单态射. 而图
\[
\begin{tikzcd}[column sep=large,row sep=large]
\Hom_{\mathbb Z}(\Hom_A(E,F'),G)
 \arrow[r]\arrow[d,"\sim"']&
\Hom_{\mathbb Z}(\Hom_A(E,F),G)\arrow[d,"\sim"]\\
\Hom_{\mathbb Z}(F',G)\otimes_A E
 \arrow[r]&
\Hom_{\mathbb Z}(F,G)\otimes_A E
\end{tikzcd}
\]
交换, 且由关于 $E$ 的假设, 纵向箭头都是同构.

\emph{命题 5.2 的证明.} 蕴含 a) $\Rightarrow$ b) 是直接的. 为证明 b) $\Rightarrow$ a), 设
\[
\cdots\longrightarrow P_n\xrightarrow{d_n}\cdots\longrightarrow P_i\longrightarrow 0\longrightarrow\cdots
\]
为在 $D(\cC)$ 中同构于 $L^\bullet$ 的复形, 其中各 $P_j$ 为有限型投射 $A$-模. 由于 $L^\bullet$ 有有限 Tor 维数, 且上同调有界, 存在 $n_o\in\mathbb Z$ 使得:

I) 序列
\[
\cdots\longrightarrow P_{n_o-p}\xrightarrow{d_{n_o-p}}\cdots\longrightarrow
P_{n_o}\longrightarrow \Ker d_{n_o+1}\longrightarrow0
\]
正合.

II) $\Ker d_{n_o+1}$ 平坦.

III) 复形
\[
0\longrightarrow \Ker d_{n_o+1}\longrightarrow P_{n_o+1}\longrightarrow\cdots\longrightarrow P_i\longrightarrow0\longrightarrow\cdots
\]
在 $D(\cC)$ 中同构于 $L^\bullet$.

而 $\Ker d_{n_o+1}$ 为有限表现, 因而根据引理为有限型投射模, 这就证明了 8.1.

若 $A$ 为交换环, $X^\bullet,Y^\bullet$ 为两个完美复形, 则其在导出范畴中的张量积
\[
X^\bullet\overset{\mathbf L}{\otimes}_A Y^\bullet
\]
仍为完美复形. 得到双函子
\[
D(A)_{\parf}\times D(A)_{\parf}\longrightarrow D(A)_{\parf}
\]
它关于每个变量都正合. 因此在 $K^\bullet(A)$ 上定义乘法, 容易看出这样便赋予 $K^\bullet(A)$ 含幺交换环结构. 实际上, 可赋予 $K^\bullet(A)$ 一个 $\lambda$-环结构; 详见继本讨论班之后的 Riemann--Roch 讨论班.\footnote{SGA 6 V 2.}

完美复形 $X^\bullet$ 与有界伪凝聚复形 $Y^\bullet$ 的张量积
\[
X^\bullet\overset{\mathbf L}{\otimes}_A Y^\bullet
\]
为有界伪凝聚复形 (仍假设 $A$ 为交换环). 因此可在 $K_\bullet(A)$ 上赋予 $K^\bullet(A)$-模结构.

\subsection*{9. 函子性质}

a) 设 $A,B$ 为两个不必交换的环, $u:A\to B$ 为同态. 每个左 (分别地右) $B$-模 $M$ 都可赋予左 (分别地右) $A$-模结构, 记为 $M_{[u]}$. 特别地, $B_{[u]}$ 为右 $A$-模.

函子
\[
E\longmapsto B_{[u]}\otimes_A E
\]
给出正合函子
\[
D(A)_{\parf}\longrightarrow D(B)_{\parf}
\]
进而给出同态
\[
K^\bullet(A)\xrightarrow{u^*}K^\bullet(B).
\]
于是得到从环 (分别地交换环) 范畴到阿贝尔群 (分别地交换环) 范畴的协变函子.

b) 现在假设左 $A$-模 $B_{[u]}$ 伪凝聚. 因而函子
\[
F\longmapsto F_{[u]}
\]
将有限型自由 $B$-模变成伪凝聚 $A$-模. 利用命题 5.1, 得到正合函子
\[
D^b(B)_{\coh}\longrightarrow D^b(A)_{\coh}
\]
进而得到典范同态
\[
K_\bullet(B)\xrightarrow{u_*}K_\bullet(A).
\]
同态 $u_*$ 和 $u^*$ 满足下列恒等式 (“投影公式”):
\[
u_*(u^*(\alpha)\cdot \beta)=\alpha\,(u_*(\beta))
\quad\text{其中 }\alpha\in K^\bullet(A),\ \beta\in K_\bullet(B).
\]
这一简单公式将在 Riemann--Roch 讨论班的报告中以更一般的形式得到证明 (SGA 6 IV 2.11.1.2).

\subsection*{10. 相对构造\footnote{参见 (SGA 6 III, IV).}}

设 $u:A\to B$ 为环同态, 并设
\[
D(B)\longrightarrow D(A)
\]
为相关联的正合函子. 在 $D(B)$ 中, 可考虑相对于 $A$ 伪凝聚的复形、相对于 $A$ 有有限 Tor 维数的复形, 以及相对于 $A$ 完美的复形. 特别地, 自然考虑 $D(B)$ 的三角子范畴
\[
D(B)_{A-\parf},\qquad D^b(B)_{A-\coh}
\]
它们分别由相对于 $A$ 完美的复形和相对于 $A$ 伪凝聚且有界的复形生成. 这两个三角范畴可用来定义下列群:
\[
K^\bullet(B\operatorname{\,rel\,}A)=K(D(B)_{A-\parf}),
\]
\[
K_\bullet(B\operatorname{\,rel\,}A)=K(D^b(B)_{A-\coh}).
\]

\emph{例.} 设 $G$ 为群, $\Lambda$ 为交换环. 同态
\[
\Lambda\xrightarrow{\varepsilon}\Lambda[G]
\]
因而可用来定义群
\[
K^\bullet(\Lambda[G]\operatorname{\,rel\,}\Lambda),\qquad
K_\bullet(\Lambda[G]\operatorname{\,rel\,}\Lambda),
\]
它们分别记为
\[
R^\bullet_\Lambda(G),\qquad R^\Lambda_\bullet(G).
\]
实际上, $R^\bullet_\Lambda(G)$ 为交换环, $R^\Lambda_\bullet(G)$ 是其上的模.

\emph{注.}

a) 若环 $\Lambda$ 不必交换, 可按定义令
\[
\Lambda[G]=\Lambda\otimes_{\mathbb Z}\mathbb ZG
\]
因而得到同态
\[
\Lambda\xrightarrow{\varepsilon}\Lambda[G].
\]
这样, 即使 $\Lambda$ 非交换, 群 $R^\bullet_\Lambda(G)$ 和 $R^\Lambda_\bullet(G)$ 也有意义.

b) 上述构造可推广到 $\Spec\Lambda$ 上的群预概形或投射群预概形 $G$ 的情形, 从带有 $G$ 作用的 $\Lambda$-模的阿贝尔范畴出发; 此时不再用环上的模明确表达. 然后使用到 $\Lambda$-模范畴的自然包含函子. 若 $\Lambda$ 为域, 可用投射环上的模明确表达; 一般情形下, 至少当 $G$ 在 $\Lambda$ 上仿射时, 可用余模明确表达; 参见 SGA 3 I 4.7.2).

\subsection*{参考文献}

\begin{enumerate}
\item[{[1]}] A. Grothendieck 与 J. Dieudonné, \emph{Elements de Géométrie Algébrique}, Publ. Math. I.H.E.S. IV (第二部分).
\item[{[2]}] J. L. Verdier, \emph{Catégories Dérivées}, I.H.E.S., 1963.\footnote{发表于 SGA $4^{1/2}$.}
\end{enumerate}
